% !TeX root = ../../brightPhD.tex
\chapter{拓扑优化模型与数值基础}
\label{chap:top_models_numerics}

%\iffalse

\section{主要符号表与基本函数空间}
\label{sec:notation_spaces}

建立严密的数学模型与数值离散算法，离不开规范的数学符号与函数空间体系。本节将首先梳理贯穿全文的主要物理量与数学符号，随后详细给出用于构建有限元变分框架的基本 Sobolev 空间及其范数定义，为后续章节中拓扑优化模型的连续化推导与有限元分析提供统一的数学基础。


\subsection{主要符号表}
\label{subsec:notation_table}

为了方便后续章节的理论表述与公式推导，本节对全文中常用的数学符号及其含义作统一约定，如表~\ref{tab:main_symbols} 所示。

\begin{table}[htbp]
	\centering
	\renewcommand{\arraystretch}{1.2} 
	\caption{本文主要数学符号及其含义}
	\label{tab:main_symbols}
	\begin{tabular}{l @{\hspace{2.5em}} p{10cm}} 
		\toprule
		符号 & 物理意义及数学说明 \\
		\midrule
		% --- 标量与参数 ---
		$d$ & 空间维数，文中通常取 $d\in\{2,3\}$ \\
		$k$ & 有限元阶次 \\
		\addlinespace
		
		% --- 空间与集合 ---
		$\Omega$ & $\mathbb{R}^d$ 中的有界开区域 \\
		$\mathbb{N}$ & 非负整数集，$\mathbb{N} = \{0,1,2,\cdots\}$ \\
		$\mathbb{R}$ & 实数域 \\
		$\mathbb{R}^d$ & $d$ 维欧氏空间 \\
		$\mathbb{M}$ & 二阶张量空间，$\mathbb{M} := \mathbb{R}^{d\times d}$ \\
		$\mathbb{S}$ & 二阶对称张量空间，$\mathbb{S} := \{\boldsymbol{\tau}\in\mathbb{M}:\boldsymbol{\tau}^T = \boldsymbol{\tau}\}$ \\
		\addlinespace
		
		% --- 物理场与向量 ---
		$\boldsymbol{u}$ & 位移场，$\boldsymbol{u}:\Omega\to\mathbb{R}^d$ \\
		$\boldsymbol{\varepsilon}(\boldsymbol{u})$ & 位移场 $\boldsymbol{u}$ 的对称应变场，$\boldsymbol{\varepsilon}(\boldsymbol{u}):\Omega\to\mathbb{S}$ \\
		$\boldsymbol{\sigma}$ & 柯西应力场$\boldsymbol{\sigma}:\Omega\to\mathbb{S}$ \\
		$\mathbb{C}$ & 四阶线弹性刚度张量 $\mathbb{C}:\mathbb{S}\to\mathbb{S}$，表示由对称应变到对称应力的线性映射\\
		$\mathcal{A}$ & 四阶线弹性柔度张量 $\mathcal{A}:\mathbb{S}\to\mathbb{S}$，表示由对称应力到对称应变的线性映射，在可压缩线弹性情形下可记为 $\mathbb{C}^{-1}$\\
		\addlinespace
		
		% --- 本构映射与特殊符号 ---
		$E$ & 杨氏模量\\
		$\nu$ & 泊松比\\
		$\lambda,\mu$ & 拉梅常数\\
		$\delta_{ij}$ & Kronecker 符号 \\
		\bottomrule
	\end{tabular}
\end{table}


\subsection{基本函数空间}
\label{subsec:function_spaces}

为了规范后续变分方程与有限元离散的推导，本文采用偏微分方程数值解与计算力学中标准的函数空间记号。设 $\Omega\subset\mathbb{R}^d$ 为有界开区域。下文中，对向量的绝对值 $\vert\cdot\vert$ 表示 Euclidean 范数，对二阶张量的绝对值 $\vert\cdot\vert$ 表示 Frobenius 范数。





%本文采用计算力学与偏微分方程数值解中标准的 Sobolev 空间记号。给定 $\mathbb{R}^d$ 中的有界开区域 $\Omega$，本文涉及的连续物理场均定义在以下常用的函数空间内：

\begin{description}
	\item[标量连续函数空间：] 
	\[
	C^0(\Omega) = \{v:\Omega\to\mathbb{R} : v~\text{在}~\Omega\text{ 上连续}\}.
	\]
	
	\item[向量值连续函数空间：]
	\[
	C^0(\Omega;\mathbb{R}^d) = \{\boldsymbol{v}:\Omega\to\mathbb{R}^d : \boldsymbol{v}(\boldsymbol{x}) = (v_1(\boldsymbol{x}),\cdots,v_d(\boldsymbol{x}))^{\top}, v_i\in C^0(\Omega)\}.
	\] 
	
	\item[标量本质有界函数空间：] 
	\[
	L^{\infty}(\Omega) = \{v:\Omega\to\mathbb{R} : \operatorname*{ess\,sup}_{\boldsymbol{x}\in\Omega}\vert v(\boldsymbol{x})\vert < \infty\},
	\]
	其范数定义为 
	\[
	\Vert v\Vert_{L^\infty(\Omega)} = \operatorname*{ess\,sup}_{\boldsymbol{x}\in\Omega}\vert v(\boldsymbol{x})\vert.
	\]
	
	\item[对称张量值本质有界函数空间：] 
	\[
	L^{\infty}(\Omega;\mathbb{S}) = \left\{\boldsymbol{\tau}:\Omega\to\mathbb{S} : \operatorname*{ess\,sup}_{\boldsymbol{x}\in\Omega}\vert\boldsymbol{\tau}(\boldsymbol{x})\vert <\infty \right\},
	\]
	其范数定义为
	\[
	\Vert\boldsymbol{\tau}\Vert_{L^\infty(\Omega;\mathbb{S})} = \operatorname*{ess\,sup}_{\boldsymbol{x}\in\Omega}\vert\boldsymbol{\tau}(\boldsymbol{x})\vert.
	\]
	
	\item[向量值与张量值平方可积空间：] 
	\[
	L^2(\Omega;\mathbb{R}^d) = \left\{\boldsymbol{v}:\Omega\to\mathbb{R}^d : \int_{\Omega} \vert\boldsymbol{v}(\boldsymbol{x})\vert^2 \,\mathrm{d}\boldsymbol{x} < \infty\right\},
	\]
	其范数定义为 
	\[
	\Vert\boldsymbol{v}\Vert_{0,\Omega} = \left(\int_{\Omega}|\boldsymbol{v}(\boldsymbol{x})|^2\,\mathrm{d}\boldsymbol{x}\right)^{1/2}.
	\]
	同理可定义张量空间 $L^2(\Omega;\mathbb{M})$ 与 $L^2(\Omega;\mathbb{S})$ 及其对应的诱导范数 $\Vert\cdot\Vert_{0,\Omega}$。
	
	\item[向量值一阶 Sobolev 空间：]
	\[
	H^1(\Omega;\mathbb{R}^d) = \{\boldsymbol{v}\in L^2(\Omega;\mathbb{R}^d) : \nabla\boldsymbol{v}\in L^2(\Omega;\mathbb{M})\},
	\] 
	其范数定义为 
	\[
	\Vert\boldsymbol{v}\Vert_{1,\Omega} = \left(\Vert\boldsymbol{v}\Vert^2_{0,\Omega} + \Vert\nabla\boldsymbol{v}\Vert^2_{0,\Omega}\right)^{1/2}.
	\]
	
	\item[散度平方可积对称张量空间：]
	\[
	H(\operatorname{div},\Omega;\mathbb{S}) = \big\{\boldsymbol{\tau}\in L^2(\Omega;\mathbb{S}) : \operatorname{div}\boldsymbol{\tau}\in L^2(\Omega;\mathbb{R}^d)\big\},
	\]
	其范数定义为 
	\[
	\Vert\boldsymbol{\tau}\Vert_{H(\operatorname{div},\Omega)} = \left(\Vert\boldsymbol{\tau}\Vert^2_{0,\Omega} + \Vert\operatorname{div}\boldsymbol{\tau}\Vert^2_{0,\Omega}\right)^{1/2}
	\]
\end{description}


\section{变密度拓扑优化的连续数学模型}
\label{sec:cont_models}

连续体拓扑优化的原始形式是一个寻找最优材料分布的无穷维 0-1 规划问题，即在给定设计域 $\Omega$ 上确定离散的材料示性函数 $\chi(\boldsymbol{x}) \in \{0, 1\}$，以最小化特定的目标泛函~\cite{bendsoeTopologyOptimization2004a}。然而，该问题在数学上通常是不适定的（即缺乏经典解的存在性），直接对其进行离散求解极易导致网格依赖性与棋盘格等严重的数值不稳定现象~\cite{kohnOptimalDesignRelaxation1986}。为克服这一理论与数值困难并有效利用基于梯度的优化算法，变密度法采用了数学松弛策略，引入连续密度变量 $\rho(\boldsymbol{x}) \in [0, 1]$ 替代离散示性变量~\cite{bendsoeOptimalShapeDesign1989a}，从而将原本棘手的组合优化转化为连续型的分布参数优化问题。本节将详细阐述基于该策略的连续数学模型，具体包括材料属性的插值方案以及一般形式的优化列式。


\subsection{材料密度与材料属性插值}
\label{subsec:material_interpolation}
在变密度拓扑优化中，材料密度与材料物理属性之间的插值模型是连接数学优化算法与物理力学分析的桥梁。我们需要建立一个从无量纲的密度场到具有物理量纲的材料张量场的映射关系。

设设计域为一个有界开集 $\Omega\subseteq\mathbb{R}^d$。引入材料密度函数 $\rho$，并要求其满足：
\[
\rho\in{L}^\infty(\Omega),\quad0\leq\rho(\boldsymbol{x})\leq1\quad \text{a.e. in } \Omega,
\]
其中，$\rho(\boldsymbol{x})=1$ 表示该点为实体材料，$\rho(\boldsymbol{x})=0$ 表示孔洞，$0<\rho(\boldsymbol{x})<1$ 表示中间密度的材料。这种允许密度值在 $[0,1]$ 区间连续变化的参数化方法，使得结构的拓扑和形状在优化过程中平滑地演变。

为了建立宏观结构性能与微观材料密度之间的联系，需要定义材料的本构属性及其对密度函数 $\rho(\boldsymbol{x})$ 的依赖关系。在线弹性小变形理论框架下，材料的应力–应变关系可由弹性刚度张量 $\mathbb{C}$ 所表征：
\[
\boldsymbol{\sigma}(\boldsymbol{u}) = \mathbb{C}(\rho):\boldsymbol{\varepsilon}(\boldsymbol{u}).
\]
材料插值模型的核心目标在于构造从标量密度场 $\rho(\boldsymbol{x})$ 到张量场 $\mathbb{C}(\rho(\boldsymbol{x}))$ 的映射
\[
\rho(\boldsymbol{x})\mapsto \mathbb{C}(\rho(\boldsymbol{x})),
\]
并要求插值得到的刚度张量保持线弹性问题适定性所需的基本性质（即对称性与一致强椭圆性）。

对于各向同性材料，弹性刚度张量 $\mathbb{C}$ 可由杨氏模量 $E$ 和泊松比 $\nu$ 这两个标量参数完全描述。在建立基本变密度模型时，此处假设泊松比 $\nu$ 在优化过程中保持为常数 $\nu_0$，而仅将杨氏模量 $E$ 视为密度 $\rho$ 的函数，即$E=E(\rho)$。在此设定下，弹性刚度张量 $\mathbb{C}$ 可直接写作 $\mathbb{C}(\rho) = \mathbb{C}(E(\rho),\nu_0)$，从而材料插值问题可等价地理解为构造标量函数 $E(\rho)$ 的问题。由于 $\rho\in L^\infty(\Omega)$ 且 $E(\cdot)$ 为定义在 $[0,1]$ 上的连续函数，因此 $E(\rho)\in L^\infty(\Omega)$。

在众多材料插值方案中，SIMP 模型是应用最为广泛的一类幂律插值 \cite{bendsoeOptimalShapeDesign1989a}。其基本思想是将等效杨氏模量 $E(\rho(\boldsymbol{x}))$ 表示为相对于实体材料杨氏模量 $E_0$ 的幂函数：
\begin{equation} \label{eq:ch6_standard_SIMP}
	E(\rho(\boldsymbol{x})) = \rho(\boldsymbol{x})^p E_0,
\end{equation}
其中 $p>1$ 为惩罚因子，惩罚因子的引入旨在显著削弱中间密度单元的刚度贡献，使得在优化过程中中间密度配置在能量上处于劣势，从而驱动最优密度场向接近 0–1 的分布收敛，以获得拓扑清晰的结构。

然而，标准 SIMP 模型在 $\rho(\boldsymbol{x})\to0$ 时给出 $E(\rho(\boldsymbol{x}))\to0$。在有限元离散中，若简单采用 $E=0$ 对应孔洞区域，则局部刚度矩阵可能退化为奇异矩阵，导致整体刚度矩阵病态甚至不可逆，进而引发数值不稳定。为避免该问题，并在物理上为孔洞区域保留极小但非零的刚度，以维持数值稳定性，工程上通常采用修正的 SIMP 模型 \cite{zhouCOCAlgorithmPart1991}，对杨氏模量设置一个正的下限 $E_{\min}>0$。修正 SIMP 插值可写为：
\begin{equation} \label{eq:ch6_modified_SIMP}
	E(\rho(\boldsymbol{x})) = E_{\min} + \rho(\boldsymbol{x})^p(E_0-E_{\min}),
\end{equation}
其中 $E_{\min}\ll{E}_0$ 为非常小的正数。这样，即使在 $\rho(\boldsymbol{x})=0$ 的区域，材料也保有微小的刚度 $E_{\min}$，从而保证整体刚度矩阵在数值上保持良好的条件数。

除 SIMP 外，文献中还提出了多种替代性材料插值模型。例如，RAMP 模型采用有理函数形式~\cite{stolpeAlternativeInterpolationScheme2001a}：
\[
E(\rho(\boldsymbol{x})) = E_{\min} + \frac{\rho(\boldsymbol{x})(E_0-E_{\min})}{1+q(1-\rho(\boldsymbol{x}))},
\]
其中 $q>0$ 为控制参数。与 SIMP 相比，RAMP 在惩罚行为、插值函数的凸性/凹性以及在低密度区域的导数特征等方面呈现不同规律，在某些情形下有利于改善优化问题的数值性质和收敛特性。

需要指出的是，尽管采用了惩罚型材料插值，在实际优化迭代过程中（尤其是中间阶段以及收敛解附近）通常仍会出现一定比例的中间密度区域，即所谓灰度区域。这些灰度区域在物理解释上对应局部材料尺度未明确的状态，不利于制造实现。同时，它们也反映出当前插值与离散框架下中间密度难以被充分抑制。若进一步缺乏适当的过滤、投影或其他正则化机制，则离散求解中还可能伴随棋盘格、网格依赖性及局部极值等典型数值问题。因此，如何在保持结构全局力学性能的同时，有效抑制过度的灰度区域、提高拓扑的清晰性与可制造性，并兼顾离散求解的数值稳定性，是变密度方法理论研究和工程应用中必须审慎处理的重要课题。


\subsection{变密度拓扑优化的一般列式}
\label{subsec:general_formulation}

设设计域为有界开集 $\Omega\subset\mathbb{R}^d$，根据物理边界条件，可将边界 $\partial\Omega$ 分解为狄利克雷边界 $\Gamma_D$ 与诺伊曼边界 $\Gamma_N$，使其满足  $\Gamma_D \cup \Gamma_N = \partial\Omega$  且 $\Gamma_D \cap \Gamma_N = \emptyset$。记状态变量为位移场 $\boldsymbol{u}:\Omega\to\mathbb{R}^d$，并要求 $\boldsymbol{u} \in \mathcal{V}$。此处 $\mathcal{V} \subset H^1(\Omega;\mathbb{R}^d)$ 表示满足给定狄利克雷边界条件的试探函数空间，同时，引入与之对应的测试函数空间 $\mathcal{V}_0$，定义为在边界 $\Gamma_D$ 上满足齐次边界条件的 $H^1$ 向量值函数集合。

设计变量为标量密度场 $\rho:\Omega\to[0,1]$，其在本质有界空间 $L^\infty(\Omega)$ 中构成的可行设计集合 $\mathcal{X}$ 的典型形式为
\[
\mathcal{X} = \{\rho(\boldsymbol{x})\in{L}^\infty(\Omega):0\leq\rho(\boldsymbol{x})\le1\quad\text{a.e.}~\text{in}\,\Omega\}.
\]

在变密度方法框架下，连续体结构拓扑优化可以抽象为以下偏微分方程约束优化问题：
\begin{equation} \label{eq:ch6_general_topology_optimization}
	\begin{aligned}
		\min_{\rho,\boldsymbol{u}}\quad&\mathcal{J}(\boldsymbol{u}, \rho)\\
		\text{subject~to}\quad&
		a_\rho(\boldsymbol{u},\boldsymbol{v}) = \ell(\boldsymbol{v}),\quad\forall\boldsymbol{v}\in\mathcal{V}_0,\\
		&\mathcal{G}_i(\rho,\boldsymbol{u})\leq0,\quad{i}=1,\cdots,m,\\
		&\rho\in\mathcal{X}.
	\end{aligned}
\end{equation}
其中各项含义说明如下：
\begin{itemize}
	\item 目标泛函 $\mathcal{J}(\boldsymbol{u}, \rho)$：用于衡量结构宏观性能的泛函，其映射关系为 $\mathcal{J}:\mathcal{V}\times\mathcal{X}\to\mathbb{R}$，具体形式随所考虑的优化问题而定。如柔顺度最小化、柔顺机构设计与体积最小化等典型目标。

	\item 状态方程 $a_\rho(\boldsymbol{u},\boldsymbol{v}) = \ell(\boldsymbol{v})$：由密度场 $\rho$ 决定的物理场平衡方程的弱形式。双线性型 $a_\rho(\cdot,\cdot)$ 通过材料刚度张量 $\mathbb{C}(\rho)$ 依赖于材料密度函数 $\rho(\boldsymbol{x})$，线性泛函 $\ell(\cdot)$ 则对应外载荷与体力等源项的作用。
	
	\item 约束条件 $\mathcal{G}_i(\rho,\boldsymbol{u})\leq 0$ ($i=1,\cdots,m$)：用于刻画设计必须满足的工程与物理要求，其中 $m$ 为不等式约束的总个数。如全局体积约束与局部应力约束等典型的不等式约束。
	
	\item 设计变量可行集合 $\mathcal{X}$：限制密度函数 $\rho(\boldsymbol{x})$ 的取值范围与基本正则性，这不仅赋予了优化结果以明确的物理意义，更是确保优化问题在数学上具备合理性质的基础。
\end{itemize}


\section{变密度拓扑优化的典型问题}
\label{sec:typical_problems}

在第~\ref{subsec:general_formulation} 节建立的偏微分方程约束优化一般框架下，本节将针对工程中经典的三类结构设计需求，给出其具体的数学列式与变分形式。这三类问题涵盖了从全局能量极值（柔顺度最小化）、特定边界响应控制（柔顺机构设计）到局部高度非线性约束（应力约束）的典型变分场景，它们不仅在数学性质与灵敏度分析策略上各具代表性，同时也为本文后续章节推导混合有限元离散格式以及设计高效非线性求解算法提供了标准的数学模型与数值验证基础。


\subsection{柔顺度最小化问题}
\label{subsec:min_compliance}

柔顺度度量了外载荷对结构所作的功，因此可以视作“柔软程度”的量化指标。对于给定的状态变量 $\boldsymbol{u}$ 与设计变量 $\rho$，其连续形式可定义为
\[
c(\boldsymbol{u},\rho) = \ell(\boldsymbol{u}) = \int_{\Omega}\boldsymbol{b}\cdot\boldsymbol{u}\,\mathrm{d}\boldsymbol{x} + \int_{\Gamma_N}\boldsymbol{g}\cdot\boldsymbol{u}\,\mathrm{d}s,
\]
其中 $\boldsymbol{b}$ 和 $\boldsymbol{g}$ 分别表示作用在设计域 $\Omega$ 上的给定体力与作用在诺伊曼边界 $\Gamma_N$ 上的给定面力。对于线弹性结构，在平衡状态下外力所做的功等于两倍的结构应变能，因此柔顺度也可以写成应变能的形式：
\[
c(\boldsymbol{u},\rho) = 2\left(\frac{1}{2}\int_{\Omega}(\mathbb{C}(\rho):\boldsymbol{\varepsilon}(\boldsymbol{u})):\boldsymbol{\varepsilon}(\boldsymbol{u})\,\mathrm{d}\boldsymbol{x}\right) = \int_{\Omega}(\mathbb{C}(\rho):\boldsymbol{\varepsilon}(\boldsymbol{u})):\boldsymbol{\varepsilon}(\boldsymbol{u})\,\mathrm{d}\boldsymbol{x},
\]
其中 $\mathbb{C}(\rho)$ 为密度相关的四阶弹性刚度张量（其具体的材料插值关系见第~\ref{subsec:material_interpolation} 节），冒号 “$:$” 表示张量之间的双点积。

设设计域的总体积为 $V_0 = \int_{\Omega} 1\,\mathrm{d}\boldsymbol{x}$，对于给定的密度场 $\rho$，其实际占用的材料体积定义为
\[
V(\rho) = \int_{\Omega}\rho(\boldsymbol{x})\,\mathrm{d}\boldsymbol{x}.
\]
给定材料的体积分数上限 $V_f\in(0,1]$，体积分数约束要求材料的实际体积与设计域总体积之比不超过给定的上限 $V_f$，即 $\frac{V(\rho)}{V_0} \leq V_f$。对应于一般优化框架中的约束项，该不等式约束可等价地记为
\[
g_V(\rho) := \frac{V(\rho)}{V_0} - V_f \leq 0.
\]

将状态方程视作一种从设计变量到位移解的映射，结合第~\ref{subsec:general_formulation} 节的一般优化框架，体积分数约束下的连续体柔顺度最小化问题可严格表述为如下的偏微分方程约束优化问题：
\begin{equation} \label{eq:ch6_min_compliance_problem}
	\begin{aligned}
		\min_{\rho,\boldsymbol{u}}\quad&c(\boldsymbol{u}, \rho) = \int_{\Omega}(\mathbb{C}(\rho):\boldsymbol{\varepsilon}(\boldsymbol{u})):\boldsymbol{\varepsilon}(\boldsymbol{u})\,\mathrm{d}\boldsymbol{x}\\
		\text{subject~to}\quad&a_\rho(\boldsymbol{u},\boldsymbol{v}) = \ell(\boldsymbol{v}),\quad\forall\boldsymbol{v}\in\mathcal{V}_0,\\
		&g_V(\rho) \leq 0,\\
		&\rho\in\mathcal{X}.
	\end{aligned}
\end{equation}
在材料插值采用 SIMP 等惩罚型模型的设定下，双线性型 $a_\rho(\cdot,\cdot)$ 对位移变量 $\boldsymbol{u}$ 是连续且一致强椭圆的，从而保证了对于任意固定的密度场 $\rho$，其对应的线弹性边值问题在位移空间上是适定的。然而，由于刚度张量 $\mathbb{C}(\rho)$ 对密度场 $\rho$ 的非线性依赖，整体优化问题在设计变量空间上通常是高度非凸的，从而容易出现局部极值。另一方面，若进一步缺乏适当的正则化机制，则离散求解中还可能伴随网格依赖性等典型数值不稳定现象。


\subsection{柔顺机构设计问题}
\label{subsec:compliant_mechanisms}

在变密度拓扑优化框架下，柔顺机构设计的目标不再是提高结构整体刚度，而是通过合理分配材料，使在给定输入激励作用下，输出端在指定方向上的位移响应尽可能大，同时保持结构的完整性与一定的刚度水平~\cite{bendsoeTopologyOptimization2004a}。为了与第 \ref{subsec:min_compliance} 节中的柔顺度最小化问题保持一致，并出于模型与推导上的简洁性考虑，本文中柔顺机构设计一律基于小变形线弹性假设，忽略几何与材料非线性效应。

设计域为有界开集 $\Omega\subset\mathbb{R}^d$，其边界 $\partial\Omega$ 可分解为互不交叠的四部分：
\[
\Gamma_D \cup \Gamma_{\mathrm{in}} \cup \Gamma_{\mathrm{out}} \cup \Gamma_N = \partial\Omega,
\]
其中，$\Gamma_{\mathrm{in}}$ 与 $\Gamma_{\mathrm{out}}$ 分别表示机构的输入和输出端口所在的边界子集，$\Gamma_N$ 则假定为不受力的自由表面。

为了避免出现 “机制解”（即在输入端施加载荷后结构整体刚度趋于零而产生不受控的大位移），并在连续模型中显式体现输入致动器和输出工件的刚度特性，柔顺机构设计中通常在输入端和输出端引入弹簧（或等效刚度元件）进行正则化。记输入端与输出端的特征方向向量分别为 $\boldsymbol{d}_{\mathrm{in}}$ 与 $\boldsymbol{d}_{\mathrm{out}}$（均为给定单位向量，用于提取沿指定方向的位移分量），并给定输入/输出弹簧刚度 $k_{\mathrm{in}}$ 和 $k_{\mathrm{out}}$，它们分别等效描述应变基致动器与外部工件的线性刚度。据此，在原线弹性双线性型 $a_\rho(\boldsymbol{u},\boldsymbol{v})$ 的基础上，定义附加的弹簧刚度双线性型：
\[
s(\boldsymbol{u},\boldsymbol{v}) := k_{\mathrm{in}}\int_{\Gamma_{\mathrm{in}}}(\boldsymbol{u}\cdot\boldsymbol{d}_{\mathrm{in}})(\boldsymbol{v}\cdot\boldsymbol{d}_{\mathrm{in}})\,\mathrm{d}s + k_{\mathrm{out}}\int_{\Gamma_{\mathrm{out}}}(\boldsymbol{u}\cdot\boldsymbol{d}_{\mathrm{out}})(\boldsymbol{v}\cdot\boldsymbol{d}_{\mathrm{out}})\,\mathrm{d}s,
\]
给定输入端的等效体力 $\boldsymbol{b}$ 和边界载荷密度 $\boldsymbol{g}_{\mathrm{in}}$，可定义输入载荷对应的线性泛函 $\ell_{\mathrm{in}}:\mathcal{V}\to\mathbb{R}$：
\[
\ell_{\mathrm{in}}(\boldsymbol{v}) := \int_{\Omega}\boldsymbol{b}\cdot\boldsymbol{v}\,\mathrm{d}\boldsymbol{x} + \int_{\Gamma_{\mathrm{in}}}\boldsymbol{g}_{\mathrm{in}}\cdot\boldsymbol{v}\,\mathrm{d}s,
\]
其中 $\ell_{\mathrm{in}}$ 描述了输入致动器对任意位移 $\boldsymbol{v}$ 所作的外功。于是，对每个给定的密度场 $\rho\in\mathcal{X}$，柔顺机构的位移场 $\boldsymbol{u}$ 需满足下述变分形式的平衡方程：
\begin{equation} 
	a_\rho(\boldsymbol{u},\boldsymbol{v}) + s(\boldsymbol{u},\boldsymbol{v}) = \ell_{\mathrm{in}}(\boldsymbol{v}),\quad \forall \boldsymbol{v}\in\mathcal{V}_0.
	\label{eq:ch6_variational_equilibrium}
\end{equation}
在材料插值 $\mathbb{C}(\rho)$ 满足一致强椭圆性、且 $k_{\mathrm{in}},k_{\mathrm{out}}>0$ 的条件下，\eqref{eq:ch6_variational_equilibrium} 式左端的双线性型对位移变量是连续且强椭圆的，从而保证了该状态方程的适定性。

柔顺机构设计的核心指标是输出端在指定方向上的位移。为此，引入输出位移泛函 $\ell_{\mathrm{out}}:\mathcal{V}\to\mathbb{R}$：
\[
\ell_{\mathrm{out}}(\boldsymbol{u}) := \int_{\Gamma_{\mathrm{out}}}\boldsymbol{d}_{\mathrm{out}}\cdot\boldsymbol{u}\,\mathrm{d}s,
\]
用于测量输出端沿方向 $\boldsymbol{d}_{\mathrm{out}}$ 的位移响应。引入与第~\ref{subsec:min_compliance} 节相同的体积分数约束函数 $g_V(\rho)\leq0$，则体积分数约束下柔顺机构设计问题的连续优化模型可表述为如下的偏微分方程约束优化问题：
\begin{equation} \label{eq:ch6_compliant_mechanism_problem}
	\begin{aligned} 
		\max_{\rho,\boldsymbol{u}}\quad & \ell_{\mathrm{out}}(\boldsymbol{u})\\[0.3em] 
		\text{subject to}\quad & a_\rho(\boldsymbol{u},\boldsymbol{v}) + s(\boldsymbol{u},\boldsymbol{v}) = \ell_{\mathrm{in}}(\boldsymbol{v}),\quad \forall \boldsymbol{v}\in\mathcal{V}_0,\\[0.2em] 
		& g_V(\rho)\le 0,\\[0.2em] 
		& \rho\in\mathcal{X}. 
	\end{aligned}
\end{equation}
需要指出的是，与柔顺度最小化问题不同，柔顺机构设计中的目标泛函与状态方程的输入载荷不一致，因此其灵敏度分析通常不再具有柔顺度最小化问题中的自伴随简化结构，离散求解时也更易诱发人工铰链等局部病态连接现象~\cite{sigmundDesignCompliantMechanisms1997a}。因此，该问题对后续灵敏度分析的推导以及数值求解算法的稳定性都提出了更为严苛的要求。


\subsection{应力约束问题}
\label{subsec:stress_constrained}

与前两节讨论的柔顺度最小化问题（第 \ref{subsec:min_compliance} 节）或柔顺机构设计问题（第 \ref{subsec:compliant_mechanisms} 节）侧重于结构的全局物理性能不同，应力约束问题主要关注结构各处的局部强度特征。其核心目标是在实现轻量化或提升结构刚度的同时，严格控制结构内部各点的应力水平不超过材料的强度极限，从而有效抑制局部应力集中现象，确保结构在给定的载荷环境下不发生强度失效。

在小变形线弹性理论框架下，对于给定的设计域 $\Omega$，设材料密度场为 $\rho\in\mathcal{X}$，相应的运动学容许位移场为 $\boldsymbol{u}\in\mathcal{V}$，结构内部任意点 $\boldsymbol{x}\in\Omega$ 处的柯西应力张量 $\boldsymbol{\sigma}(\boldsymbol{u})(\boldsymbol{x})$ 由广义胡克定律给出：
\[
\boldsymbol{\sigma}(\boldsymbol{u})(\boldsymbol{x}) = \mathbb{C}(\rho(\boldsymbol{x})):\boldsymbol{\varepsilon}(\boldsymbol{u})(\boldsymbol{x}), \quad \forall\boldsymbol{x}\in\Omega,
\]
其中，$\boldsymbol{\varepsilon}(\boldsymbol{u})$ 为线性应变张量，$\mathbb{C}(\rho)$ 为依赖于密度的四阶材料弹性刚度张量。

对于各向同性材料，von Mises 等效应力通常被用作判断材料是否屈服的强度准则。利用 Voigt 记号，记 $\hat{\boldsymbol{\sigma}}$ 为应力张量的向量形式，则 von Mises 应力 $\sigma^{\text{vM}}$ 可表述为二次型的平方根形式：
\[
\sigma^{\text{vM}}(\boldsymbol{u}_\rho)(\boldsymbol{x}) = \sqrt{ \hat{\boldsymbol{\sigma}}(\boldsymbol{u}_\rho)(\boldsymbol{x})^\top \boldsymbol{M} \, \hat{\boldsymbol{\sigma}}(\boldsymbol{u}_\rho)(\boldsymbol{x})},
\]
其中 $\boldsymbol{M}$ 为与应力状态相关的常数权重矩阵。在三维情形下，应力向量
\[
\hat{\boldsymbol{\sigma}} = [\sigma_{xx}, \sigma_{yy}, \sigma_{zz}, \tau_{xy}, \tau_{yz}, \tau_{zx}]^\top,
\]
对应的矩阵 $\boldsymbol{M}$ 为
\[
\boldsymbol{M} = 
\begin{bmatrix} 
	1 & -0.5 & -0.5 & 0 & 0 & 0 \\ 
	-0.5 & 1 & -0.5 & 0 & 0 & 0 \\ 
	-0.5 & -0.5 & 1 & 0 & 0 & 0 \\ 
	0 & 0 & 0 & 3 & 0 & 0 \\ 
	0 & 0 & 0 & 0 & 3 & 0 \\ 
	0 & 0 & 0 & 0 & 0 & 3 
\end{bmatrix}.
\]
在二维平面应力假设下（$\sigma_{zz} = 0$），应力向量取 $\hat{\boldsymbol{\sigma}} = [\sigma_{xx}, \sigma_{yy}, \tau_{xy}]^\top$，矩阵简化为
\[
\boldsymbol{M} = 
\begin{bmatrix} 
	1 & -0.5 & 0 \\ -0.5 & 1 & 0 \\ 0 & 0 & 3 
\end{bmatrix},
\]
而在二维平面应变假设下（$\varepsilon_{zz} = 0$），矩阵 $\boldsymbol{M}$ 与泊松比 $\nu$ 相关，可表示为
\[
\boldsymbol{M} = 
\begin{bmatrix} 
	1-\nu+\nu^2 & -(0.5+\nu-\nu^2) & 0 \\ -(0.5+\nu-\nu^2) & 1-\nu+\nu^2 & 0 \\ 0 & 0 & 3 
\end{bmatrix}.
\]

应力作为材料内部抵抗变形的力学响应，其存在以物质的存在为前提。因此，强度失效准则仅对实体材料区域以及存在部分材料的中间密度区域具有明确的物理意义，而在完全无材料支撑的理想孔洞区域中，对应力响应施加约束通常并无实际意义。记 $\Omega_{\text{void}} = \{ \boldsymbol{x} \in \Omega \mid \rho(\boldsymbol{x}) = 0 \}$ 为理想意义下的完全孔洞区域，则物理上合理的局部应力约束原则上应仅定义在 $\Omega \setminus \Omega_{\text{void}}$ 上。

基于上述定义与分析，不失一般性，包含局部应力约束的通用拓扑优化数学模型可表述为如下的偏微分方程约束优化问题：
\begin{equation} 
	\begin{aligned}
		\min_{\rho,\boldsymbol{u}}\quad&\mathcal{J}(\boldsymbol{u}, \rho)\\
		\text{subject~to}\quad&
		a_\rho(\boldsymbol{u},\boldsymbol{v}) = \ell(\boldsymbol{v}),\quad\forall\boldsymbol{v}\in\mathcal{V}_0,\\
		&\sigma^{\text{vM}}(\boldsymbol{u})(\boldsymbol{x}) \le \bar{\sigma}, \quad \forall \boldsymbol{x} \in \Omega \setminus \Omega_{\text{void}},\\
		&\rho\in\mathcal{X}.
	\end{aligned}
	\label{eq:ch6_stress_constrained_problem}
\end{equation}
其中，目标泛函 $\mathcal{J}$ 根据具体问题选取（如体积最小化），而不等式 $\sigma^{\text{vM}}(\boldsymbol{u})(\boldsymbol{x}) \le \bar{\sigma}$ 则代表了必须在非孔洞设计域内逐点满足的局部应力约束条件，$\bar{\sigma}$ 为材料的许用应力限值。

需要指出的是，\eqref{eq:ch6_stress_constrained_problem}式 所示的连续模型在离散化求解时存在显著的数值困难，即应力奇异性。在变密度法中，当设计变量 $\rho \to 0$ 时，单元刚度矩阵趋于零，但其节点位移通常为非零有限值，导致计算出的应变和应力并不为零。若直接应用 $\sigma^{\text{vM}} \le \bar{\sigma}$ 约束，优化算法将无法完全移除材料（即无法形成孔洞），因为一旦密度降为零，应力约束即被判定为违反。针对这一问题，目前文献中主要发展了两类松弛策略：

\begin{enumerate}
	\item \textbf{$\varepsilon$-松弛技术}：由 Cheng 和 Guo 等~\cite{chengErelaxedApproachStructural1997} 提出，通过引入一个依赖于密度的微小松弛参数来修正约束不等式的界限。该技术使得当密度趋于零时，许用应力界限趋于无穷大，从而在孔洞区域自然满足约束条件。
	
	\item \textbf{应力惩罚和消失约束技术}：应力惩罚由 Bruggi~\cite{bruggiAlternativeApproachStress2008a} 及 Holmberg~\cite{holmbergStressConstrainedTopology2013a} 等发展，其中 Bruggi 和 Duysinx 将刚度惩罚指数与应力惩罚指数分离，
	形成了 $q\text{-}p$ 松弛框架~\cite{bruggiTopologyOptimizationMinimum2012}。该策略通过引入基于密度的幂函数插值，强制计算应力值随密度降低而衰减至零，从而消除奇异性。消失约束最早由 Cheng 和 Jiang~\cite{chengStudyTopologyOptimization1992} 提出，采用线性松弛形式使得约束在孔洞处自动失效。在此基础上，Giraldo-Londoño 和 Paulino~\cite{giraldo-londonoUnifiedApproachTopology2020} 提出了多项式消失约束。该方法在保留传统消失约束特性的同时，引入了应力违约量的三次项，对高应力违约区域提供了更强的惩罚梯度，从而显著加速了非线性优化的收敛过程。
\end{enumerate}

除了奇异性问题外，应力约束拓扑优化在数值求解时还面临着局部约束数量庞大的挑战。在有限元离散模型中，应力评估通常是在每个单元的积分点或中心点进行的。对于网格规模较大的精细结构，若直接对每一个评估点的应力施加约束（即 $\sigma_{j} \le \bar{\sigma}, \forall j$），将导致约束函数的数量达到数万甚至百万量级。如此庞大的非线性约束集合不仅会极大地增加灵敏度分析的计算存储负担，还会严重削弱基于梯度的优化算法的收敛性能。针对这一计算规模难题，文献中主要提出了以下三种处理策略：

\begin{enumerate}
	\item \textbf{有效集策略}：由 Bruggi 和 Duysinx~\cite{bruggiTopologyOptimizationMinimum2012} 等提出，该策略保留局部约束的离散形式，但在每次迭代中仅选取当前违反约束或接近约束边界的应力评估点构建“有效约束集”参与优化求解。该方法理论上能实现对应力的精确控制，但在拓扑演化初期，有效集的频繁变动极易导致优化算法产生震荡，增加收敛难度。
	
	\item \textbf{聚合与聚类策略}：全局聚合由 Le 等~\cite{leStressbasedTopologyOptimization2010b} 系统提出，
	利用 P-范数或 Kreisselmeier-Steinhauser (KS) 函数将全域内的局部应力聚合为单个全局积分型约束。该方法计算成本最低，但由于积分的平均化效应，往往会低估实际的最大应力峰值，导致局部应力超限。分块聚类由 París~\cite{parisBlockAggregationStress2010}、Le~\cite{leStressbasedTopologyOptimization2010b} 及 Holmberg 等~\cite{holmbergStressConstrainedTopology2013a} 发展，为了缓解全局聚合的精度问题，该策略将应力评估点划分为若干个聚类组（基于几何位置或应力幅值），并在组内分别应用聚合函数。这种方法在控制局部应力精度与计算效率之间取得了较好的平衡。
	
	\item \textbf{增广拉格朗日方法 (Augmented Lagrangian Method, ALM)}：这是一种区别于约束聚合的数学规划方法，由 Pereira~\cite{pereiraTopologyOptimizationContinuum2004} 及 da Silva~\cite{dasilvaTopologyOptimizationContinuum2018} 等引入拓扑优化领域。不同于聚合策略，该方法并不通过减少约束数量来降低规模，而是通过引入拉格朗日乘子项和二次罚函数项，将大量局部应力约束吸收到目标函数中，从而将原始的有约束优化问题转化为一系列无约束子问题序列进行求解。该方法能够严格保持应力约束的局部特性，避免了聚合函数带来的精度损失，适合处理包含大规模局部约束的拓扑优化问题。
\end{enumerate}

除了上述数值困难，由理想集中载荷引发的人工应力奇异性是应力约束拓扑优化面临的另一大挑战。在线弹性有限元框架下，点载荷会使其作用点邻域的应力随网格细化而发散。这类无穷大的应力峰值将导致局部约束始终被违背，迫使优化算法在奇异点附近盲目堆积材料，造成求解发散或生成异常构型。针对该问题，现有文献通常采用两类处理策略~\cite{holmbergStressConstrainedTopology2013a, giraldo-londonoPolyStressMatlabImplementation2021}：

\begin{enumerate}
	\item \textbf{局部应力豁免与实体保留}：在载荷点附近预设一小块非设计域（被动单元），将其密度固定为 $\rho=1$；同时在应力约束构造时，将该邻域单元的应力评估点从约束集合中剔除，以避免点载荷诱发的非物理峰值长期主导可行性与迭代演化。
	
	\item \textbf{载荷分布化策略}：基于圣维南原理，将理想点载荷替换为作用于有限尺度局部边界上的等效分布牵引力，具体做法是在载荷作用点邻域选取一段边界，构造等效分布面力，使其在静力意义下与原集中点力等效，从源头上削弱点载荷诱发的人工应力奇异性。
\end{enumerate}


\section{灵敏度分析理论与方法}
\label{sec:sensitivity_analysis}

在拓扑优化的数值求解中，灵敏度分析是连接连续数学模型与梯度优化算法的核心环节。本节将在前述典型模型的基础上，首先构建基于伴随方法的一般变分框架，随后推导柔顺度、柔顺机构及应力约束等具体问题的灵敏度表达式，为后续引入高阶与混合有限元等高效离散算法奠定理论基础。


\subsection{灵敏度分析一般框架}
\label{subsec:sensitivity_framework}

在使用基于梯度的优化算法求解变密度拓扑优化问题时，核心步骤是计算目标泛函和约束泛函对设计变量的灵敏度。灵敏度刻画了当设计变量发生微小改变时，目标或约束的变化率，从而为寻找最优设计提供方向。本节阐述一种计算灵敏度的通用高效方法，即伴随方法。为表述清晰，首先以单个目标泛函为例展开推导，其结论可直接推广至各类约束泛函。

根据第 \ref{subsec:general_formulation} 节的描述，对每个给定的 $\rho$，令 $\boldsymbol{u}_\rho$ 为相应状态方程的解，则目标泛函可以写为 $\mathcal{J}(\boldsymbol{u}_\rho, \rho)$。若直接根据链式法则计算 $\mathcal{J}(\boldsymbol{u}_\rho, \rho)$ 关于 $\rho$ 的导数，将不可避免地涉及状态解 $\boldsymbol{u}_\rho$ 对设计变量的导数 $\partial\boldsymbol{u}_\rho/\partial\rho$。在有限元空间离散后，该项对应于一个维度极高的稠密雅可比矩阵，显式构造和存储的计算代价极为高昂。伴随方法的核心思想，是通过构造适当的拉格朗日泛函 $\mathcal{L}$ 并引入伴随变量 $\boldsymbol{\lambda}$，从而在解析层面消除对 $\partial\boldsymbol{u}_\rho/\partial\rho$ 的显式依赖，最终以与设计变量维数基本无关的代价获得完整的梯度信息。

基于泛函分析中的拉格朗日乘子理论~\cite{luenbergerOptimizationVectorSpace1969}，引入拉格朗日泛函
\[
\mathcal{L}(\rho,\boldsymbol{u},\boldsymbol{\lambda}) := \mathcal{J}(\boldsymbol{u}, \rho) + \mathcal{R}(\rho,\boldsymbol{u})(\boldsymbol\lambda),
\]
其中残差算子定义为
\begin{equation}
	\mathcal{R}(\rho,\boldsymbol{u})(\boldsymbol{v}) := a_\rho(\boldsymbol{u},\boldsymbol{v}) - \ell(\boldsymbol{v}),\quad\forall\boldsymbol{v}\in\mathcal{V}_0.
	\label{eq:ch2_residual_operator}
\end{equation}
当 $\boldsymbol{u}=\boldsymbol{u}_\rho$ 满足状态方程时，有
\[
\mathcal{R}(\rho,\boldsymbol{u}_\rho)(\boldsymbol{v}) = 0,\quad\forall\boldsymbol{v}\in\mathcal{V}_0,
\]
从而对任意伴随变量 $\boldsymbol{\lambda}\in\mathcal{V}_0$ 都有 $\mathcal{R}(\rho,\boldsymbol{u}_\rho)(\boldsymbol{\lambda}) = 0$。因此，约化目标泛函 $\tilde{\mathcal{J}}(\rho) := \mathcal{J}(\boldsymbol{u}_\rho, \rho)$ 沿任意容许扰动方向 $\delta\rho$ 的一阶变分，等于拉格朗日泛函 $\mathcal{L}$ 在点 $(\rho,\boldsymbol{u}_\rho,\boldsymbol{\lambda})$ 处关于 $\rho$ 的一阶变分，即
\[
\delta\tilde{\mathcal{J}}(\rho) = \delta\mathcal{L}(\rho,\boldsymbol{u}_\rho,\boldsymbol{\lambda}),
\]
该一阶变分本质上给出了目标泛函关于设计变量 $\rho$ 的 Gâteaux 导数。

对拉格朗日泛函 $\mathcal{L}(\rho,\boldsymbol{u},\boldsymbol{\lambda})$ 作一阶变分，可写为各项偏 Gâteaux 导数的线性叠加~\cite{luenbergerOptimizationVectorSpace1969}：
\begin{equation} 
	\delta\mathcal{L}(\rho,\boldsymbol{u},\boldsymbol{\lambda}) = \frac{\partial\mathcal{L}(\rho,\boldsymbol{u},\boldsymbol{\lambda})}{\partial\rho}[\delta\rho] + \frac{\partial\mathcal{L}(\rho,\boldsymbol{u},\boldsymbol{\lambda})}{\partial\boldsymbol{u}}[\delta\boldsymbol{u}] + \frac{\partial\mathcal{L}(\rho,\boldsymbol{u},\boldsymbol{\lambda})}{\partial\boldsymbol\lambda}[\delta\boldsymbol\lambda].
	\label{eq:ch2_lagrangian_variation}
\end{equation}
结合 $\mathcal{L}$ 与 $\mathcal{R}$ 的定义，\eqref{eq:ch2_lagrangian_variation} 式各项可分别具体展开为
\[
\frac{\partial\mathcal{L}}{\partial\rho}[\delta\rho] = \frac{\partial\mathcal{J}( \boldsymbol{u},\rho)}{\partial\rho}[\delta\rho] + \frac{\partial\mathcal{R}(\rho,\boldsymbol{u})(\boldsymbol\lambda)}{\partial\rho}[\delta\rho],
\]
\[
\frac{\partial\mathcal{L}}{\partial\boldsymbol{u}}[\delta\boldsymbol{u}] = \frac{\partial\mathcal{J}( \boldsymbol{u},\rho)}{\partial\boldsymbol{u}}[\delta\boldsymbol{u}] + \frac{\partial\mathcal{R}(\rho,\boldsymbol{u})(\boldsymbol\lambda)}{\partial\boldsymbol{u}}[\delta\boldsymbol{u}],
\]
\[
\frac{\partial\mathcal{L}}{\partial\boldsymbol\lambda}[\delta\boldsymbol\lambda] = \mathcal{R}(\rho,\boldsymbol{u})(\delta\boldsymbol{\lambda}) = a_\rho(\boldsymbol{u},\delta\boldsymbol\lambda) - \ell(\delta\boldsymbol\lambda).
\]
特别地，当 $\boldsymbol{u}=\boldsymbol{u}_\rho$ 满足状态方程时，残差恒为零，即 $\mathcal{R}(\rho,\boldsymbol{u}_\rho)(\delta\boldsymbol{\lambda})=0$，从而有 $\frac{\partial\mathcal{L}(\rho,\boldsymbol{u}_\rho,\boldsymbol{\lambda})}{\partial\boldsymbol\lambda}[\delta\boldsymbol\lambda] = 0$。

为了消除目标泛函变分中对未知状态扰动 $\delta\boldsymbol{u}$ 的隐式依赖，引入伴随变量 $\boldsymbol{\lambda}\in\mathcal{V}_0$，并强制要求拉格朗日泛函对状态变量的偏变分为零：
\[
\frac{\partial\mathcal{L}}{\partial\boldsymbol{u}}[\delta\boldsymbol{u}] = 0, \quad \forall\,\delta\boldsymbol{u}\in\mathcal{V}_0,
\]
由拉格朗日泛函的定义展开可得
\[
\frac{\partial\mathcal{L}}{\partial\boldsymbol{u}}[\delta\boldsymbol{u}] = \frac{\partial\mathcal{J}( \boldsymbol{u}_\rho,\rho)}{\partial\boldsymbol{u}}[\delta\boldsymbol{u}] + \frac{\partial\mathcal{R}(\rho,\boldsymbol{u}_\rho)(\boldsymbol\lambda)}{\partial\boldsymbol{u}}[\delta\boldsymbol{u}],
\]
在变密度线弹性力学框架中，残差算子定义为 $\mathcal{R}(\rho,\boldsymbol{u})(\boldsymbol{v}) := a_\rho(\boldsymbol{u},\boldsymbol{v}) - \ell(\boldsymbol{v})$，因此
\[
\frac{\partial\mathcal{R}(\rho,\boldsymbol{u}_\rho)(\boldsymbol\lambda)}{\partial\boldsymbol{u}}[\delta\boldsymbol{u}] = a_\rho(\delta\boldsymbol{u},\boldsymbol{\lambda}),
\]
于是，通过将任意扰动 $\delta\boldsymbol{u}$ 替换为测试空间中的任意测试函数 $\boldsymbol{v}$，伴随变量 $\boldsymbol{\lambda}$ 应满足如下的伴随方程：求 $\boldsymbol{\lambda}\in\mathcal{V}_0$，使得
\begin{equation}
	a_\rho(\boldsymbol{v},\boldsymbol{\lambda}) = -\frac{\partial\mathcal{J}( \boldsymbol{u}_\rho,\rho)}{\partial\boldsymbol{u}}[\boldsymbol{v}],\quad \forall\,\boldsymbol{v}\in\mathcal{V}_0
	\label{eq:ch2_adjoint_equation}
\end{equation}
一旦对给定的设计变量 $\rho$ 求解状态方程得到 $\boldsymbol{u}_\rho$​，再求解 \eqref{eq:ch2_adjoint_equation} 式所示的伴随方程得到对应的伴随解 $\boldsymbol{\lambda}$，则约化目标泛函 $\tilde{\mathcal{J}}(\rho)$ 沿方向 $\delta\rho$ 的一阶变分即可大幅简化为：
\begin{equation} 
	\delta\tilde{\mathcal{J}}(\rho) = \delta\mathcal{L}(\rho, \boldsymbol{u}\rho, \boldsymbol{\lambda}) = \frac{\partial\mathcal{L}}{\partial\rho}[\delta\rho] = \frac{\partial\mathcal{J}(\boldsymbol{u}\rho, \rho)}{\partial\rho}[\delta\rho] + \frac{\partial\mathcal{R}(\rho, \boldsymbol{u}\rho)(\boldsymbol{\lambda})}{\partial\rho}[\delta\rho].
	\label{eq:ch2_gateaux_derivative}
\end{equation}
该表达式即为目标泛函关于设计变量 $\rho$ 的精确 Gâteaux 导数。此时，导数公式中已经彻底剥离了难以计算的未知状态扰动 $\delta\boldsymbol{u}$，而仅依赖于给定的状态解 $\boldsymbol{u}_\rho$、伴随解 $\boldsymbol{\lambda}$ 以及物理算子对 $\rho$ 的显式偏导数。

在变密度拓扑优化中，\eqref{eq:ch2_gateaux_derivative} 式所示的线性泛函可进一步写成设计域上的积分形式：
\[
\delta\tilde{\mathcal{J}}(\rho) = \int_{\Omega} \omega(\boldsymbol{x}) \delta\rho(\boldsymbol{x}) \,\mathrm{d}\boldsymbol{x},
\]
其中积分核 $\omega(\boldsymbol{x})$ 刻画了局部密度的微小扰动对全局目标的影响，通常称为灵敏度密度函数。伴随方法的核心优势在于其计算复杂度与设计变量维数彻底解耦。无论空间离散后的设计变量（如单元或节点密度）规模多大，获取全场灵敏度的主要计算代价仅为求解一次状态方程和一次伴随方程。这一高效的求解性质，使其成为大规模拓扑优化问题中计算梯度的标准范式。


\subsection{典型拓扑优化问题的灵敏度分析}
\label{subsec:sensitivity_derivation}

\subsubsection*{柔顺度目标函数的灵敏度}

考虑体积分数约束下的柔顺度最小化问题，对给定的密度场 $\rho\in\mathcal{X}$，记 $\boldsymbol{u}_\rho\in\mathcal{V}$ 为线弹性平衡方程弱形式的解，即满足
\begin{equation}
	a_\rho(\boldsymbol{u}_\rho,\boldsymbol{v}) = \ell(\boldsymbol{v}),\quad \forall \boldsymbol{v}\in\mathcal{V}_0.
	\label{eq:ch2_state_equation}
\end{equation}
在平衡状态下，柔顺度定义为外力所作的功，可记为关于密度的约化泛函
\[
c(\rho) = \tilde{\mathcal{J}}(\rho) = \mathcal{J}(\boldsymbol{u}_\rho,\rho) := \ell(\boldsymbol{u}_\rho),
\]
于是，目标泛函 $\mathcal{J}(\boldsymbol{u},\rho)$ 对位移 $\boldsymbol{u}$ 的偏 Gâteaux 导数为
\[
\frac{\partial\mathcal{J}(\boldsymbol{u}_\rho,\rho)}{\partial\boldsymbol{u}}[\boldsymbol{v}] = \ell(\boldsymbol{v}),\quad \forall \boldsymbol{v}\in\mathcal{V}_0.
\]
且该泛函对密度 $\rho$ 没有显式依赖，即
\[
\frac{\partial\mathcal{J}(\boldsymbol{u}_\rho,\rho)}{\partial\rho}[\delta\rho] = 0.
\]
根据第~\ref{subsec:sensitivity_framework} 节给出的伴随方法一般框架，对应的伴随方程为求解 $\boldsymbol{\lambda}\in\mathcal{V}_0$，使得
\[
a_\rho(\boldsymbol{v},\boldsymbol{\lambda}) = -\frac{\partial\mathcal{J}( \boldsymbol{u}_\rho,\rho)}{\partial\boldsymbol{u}}[\boldsymbol{v}],\quad \forall\,\boldsymbol{v}\in\mathcal{V}_0.
\]
将外力功泛函对位移的导数代入上式，并利用原状态方程~\eqref{eq:ch2_state_equation}，可得
\[
a_\rho(\boldsymbol{v},\boldsymbol{\lambda}) = -\ell(\boldsymbol{v}) = -a_\rho(\boldsymbol{u}_\rho,\boldsymbol{v}), \quad \forall\,\boldsymbol{v}\in\mathcal{V}_0.
\]
进一步，利用双线性型 $a_\rho(\cdot,\cdot)$ 的对称性，上式可等价改写为
\[
a_\rho(\boldsymbol{v},\boldsymbol{\lambda}) = a_\rho(\boldsymbol{v},-\boldsymbol{u}_\rho), \quad \forall\,\boldsymbol{v}\in\mathcal{V}_0.
\]
基于双线性型的强椭圆性（保证了变分问题解的唯一性），通过对比等式两端即可严格得出伴随解
\[
\boldsymbol{\lambda} = -\boldsymbol{u}_\rho.
\]
这表明，对于柔顺度最小化问题，伴随场与位移场仅相差一个负号，无需额外求解新的边值问题，从而进一步降低了灵敏度分析的代价。这种伴随解与状态解的强关联性，正是柔顺度最小问题自伴随性质的体现~\cite{bendsoeTopologyOptimization2004a}。

下面利用一般形式的灵敏度表达式计算柔顺度目标的密度导数。根据残差算子~\eqref{eq:ch2_residual_operator} 的定义可知其关于设计变量 $\rho$沿方向 $\delta\rho$ 的偏 Gâteaux 导数为
\[
\frac{\partial\mathcal{R}(\rho,\boldsymbol{u}_\rho)(\boldsymbol{\lambda})}{\partial\rho}[\delta\rho] = \frac{\partial a_\rho(\boldsymbol{u}_\rho,\boldsymbol{\lambda})}{\partial\rho}[\delta\rho] = \int_{\Omega}\big(\mathbb{C}'(\rho)[\delta\rho]: \boldsymbol{\varepsilon}(\boldsymbol{u}_\rho)\big):\boldsymbol{\varepsilon}(\boldsymbol{\lambda})\,\mathrm{d}\boldsymbol{x},
\]
其中 $\mathbb{C}'(\rho)[\delta\rho]$ 表示弹性刚度张量对密度的方向导数。注意到在空间每一点 $\boldsymbol{x}\in\Omega$ 处，该方向导数关于密度扰动 $\delta\rho(\boldsymbol{x})$ 是局部的且呈线性关系，因此可写为
\begin{equation}
	\mathbb{C}'(\rho)[\delta\rho](\boldsymbol{x}) = \mathbb{C}'(\rho(\boldsymbol{x}))\,\delta\rho(\boldsymbol{x}),
	\label{eq:ch2_local_linear_relationship}
\end{equation}
将其代入变分式，并利用已求得的自伴随解 $\boldsymbol{\lambda} = -\boldsymbol{u}_\rho$，可得
\[
\frac{\partial\mathcal{R}(\rho,\boldsymbol{u}_\rho)(\boldsymbol{\lambda})}{\partial\rho}[\delta\rho] = -\int_{\Omega}\big(\mathbb{C}'(\rho): \boldsymbol{\varepsilon}(\boldsymbol{u}_\rho)\big):\boldsymbol{\varepsilon}(\boldsymbol{u}_\rho)\,\delta\rho\,\mathrm{d}\boldsymbol{x}.
\]
另一方面，由于柔顺度目标泛函对 $\rho$ 无显式依赖（即 $\partial\mathcal{J}(\boldsymbol{u}_\rho,\rho)/\partial\rho[\delta\rho]=0$），根据第~\ref{subsec:sensitivity_framework} 节的一般伴随理论的结论，柔顺度约化泛函 $\tilde{\mathcal{J}}(\rho)$ 沿方向 $\delta\rho$ 的一阶总变分为
\[
\delta c(\rho) = \tilde{\mathcal{J}}(\rho) = \frac{\partial\mathcal{R}(\rho,\boldsymbol{u}_\rho)(\boldsymbol{\lambda})}{\partial\rho}[\delta\rho] = -\int_{\Omega}\big(\mathbb{C}'(\rho): \boldsymbol{\varepsilon}(\boldsymbol{u}_\rho)\big):\boldsymbol{\varepsilon}(\boldsymbol{u}_\rho)\,\delta\rho\,\mathrm{d}\boldsymbol{x},
\]
将其与变分的标准内积积分形式
\[
\delta c(\rho) = \int_{\Omega}\frac{\delta c}{\delta\rho}(\boldsymbol{x})\,\delta\rho(\boldsymbol{x})\,\mathrm{d}\boldsymbol{x},
\]
进行对比，即可直接提取出柔顺度关于密度场的灵敏度密度函数：
\begin{equation} \label{eq:ch2_compliance_sensitivity_general}
	\frac{\delta c}{\delta\rho}(\boldsymbol{x}) = -\big(\mathbb{C}'(\rho(\boldsymbol{x})): \boldsymbol{\varepsilon}(\boldsymbol{u}_\rho(\boldsymbol{x}))\big): \boldsymbol{\varepsilon}(\boldsymbol{u}_\rho(\boldsymbol{x})).
\end{equation}

若采用最经典的 SIMP 材料插值模型，弹性刚度张量及其导数可写为
\[
\mathbb{C}(\rho(\boldsymbol{x})) = E(\rho(\boldsymbol{x}))\,\mathbb{C}_0,\quad\mathbb{C}'(\rho(\boldsymbol{x})) = E'(\rho(\boldsymbol{x}))\,\mathbb{C}_0,
\]
则有
\[
\frac{\delta c}{\delta\rho}(\boldsymbol{x}) = -E'(\rho(\boldsymbol{x}))\, \big(\mathbb{C}_0:\boldsymbol{\varepsilon}(\boldsymbol{u}_\rho(\boldsymbol{x}))\big): \boldsymbol{\varepsilon}(\boldsymbol{u}_\rho(\boldsymbol{x})).
\]


\subsubsection*{柔顺机构输出位移目标的灵敏度}

考虑第~\ref{subsec:compliant_mechanisms} 节中的柔顺机构设计问题。对于给定的密度场 $\rho\in\mathcal{X}$，其位移场 $\boldsymbol{u}_\rho\in\mathcal{V}$ 由下述包含附加弹簧刚度的变分问题唯一确定：
\[
a_\rho(\boldsymbol{u}_\rho,\boldsymbol{v}) + s(\boldsymbol{u}_\rho,\boldsymbol{v}) = \ell_{\mathrm{in}}(\boldsymbol{v}),\quad \forall\,\boldsymbol{v}\in\mathcal{V}_0.
\]
柔顺机构设计的性能指标取为输出端沿指定特征方向 $\boldsymbol{d}_{\mathrm{out}}$ 的位移，这可由一个连续的线性泛函 $\ell_{\mathrm{out}}(\boldsymbol{u})$ 测量得出
\[
\ell_{\mathrm{out}}(\boldsymbol{u}) := \int_{\Gamma_{\mathrm{out}}}\boldsymbol{d}_{\mathrm{out}}\cdot\boldsymbol{u}\,\mathrm{d}s.
\]
于是，对应的一般目标泛函可定义为 $\mathcal{J}(\boldsymbol{u}, \rho) := \ell_{\mathrm{out}}(\boldsymbol{u})$。当强制满足系统状态方程时，该目标转化为仅依赖于设计变量密度的约化泛函
\[
u_{\mathrm{out}}(\rho) = \tilde{\mathcal{J}}(\rho) := \mathcal{J}(\boldsymbol{u}_\rho,\rho) = \ell_{\mathrm{out}}(\boldsymbol{u}_\rho).
\]
由于在连续模型中，一般目标泛函 $\mathcal{J}$ 对密度场 $\rho$ 没有显式依赖，且线性泛函 $\ell_{\mathrm{out}}(\cdot)$ 的导数即为其自身，因此各项偏 Gâteaux 导数可直接写为
\[
\frac{\partial\mathcal{J}(\boldsymbol{u}_\rho,\rho)}{\partial\rho}[\delta\rho] = 0,\quad\frac{\partial\mathcal{J}(\boldsymbol{u}_\rho,\rho)}{\partial\boldsymbol{u}}[\boldsymbol{v}] = \ell_{\mathrm{out}}(\boldsymbol{v}), \qquad \forall\,\boldsymbol{v}\in\mathcal{V}_0.
\]

按照一般伴随框架，将柔顺机构的状态方程写成残差算子
\[
\mathcal{R}(\rho,\boldsymbol{u})(\boldsymbol{v}) := a_\rho(\boldsymbol{u},\boldsymbol{v}) + s(\boldsymbol{u},\boldsymbol{v}) - \ell_{\mathrm{in}}(\boldsymbol{v}), \qquad \forall\,\boldsymbol{v}\in\mathcal{V}_0,
\]
并构造相应的拉格朗日泛函
\[
\mathcal{L}(\rho,\boldsymbol{u},\boldsymbol{\lambda}) := \mathcal{J}(\boldsymbol{u},\rho)  + \mathcal{R}(\rho,\boldsymbol{u})(\boldsymbol{\lambda}),\qquad\boldsymbol{\lambda}\in\mathcal{V}_0.
\]
引入伴随变量 $\boldsymbol{\lambda}$，并强制要求拉格朗日泛函对状态变量的偏变分为零：
\[
\frac{\partial\mathcal{L}(\rho,\boldsymbol{u}_\rho,\boldsymbol{\lambda})} {\partial\boldsymbol{u}}[\delta\boldsymbol{u}] = 0, \qquad \forall\,\delta\boldsymbol{u}\in\mathcal{V}_0.
\]
将前述偏导数关系代入，并将任意扰动方向 $\delta\boldsymbol{u}$ 记作测试空间中的函数 $\boldsymbol{v}$，可导出相应的伴随问题：求 $\boldsymbol{\lambda}\in\mathcal{V}_0$，使得
\[
a_\rho(\boldsymbol{v},\boldsymbol{\lambda}) + s(\boldsymbol{v},\boldsymbol{\lambda}) = -\,\ell_{\mathrm{out}}(\boldsymbol{v}), \qquad \forall\,\boldsymbol{v}\in\mathcal{V}_0.
\]
可见，伴随方程与原状态方程在左端算子形式上完全相同，但其右端不再是实际输入载荷 $\ell_{\mathrm{in}}$，而是与输出位移目标相对应的“虚拟载荷” $-\ell_{\mathrm{out}}$。因此，与柔顺度最小化问题不同，柔顺机构设计中的目标泛函与状态方程输入载荷并不一致，其灵敏度分析一般不具有柔顺度最小化问题中的自伴随简化结构，通常不能得到 $\boldsymbol{\lambda}=\pm\boldsymbol{u}_\rho$ 这类简单关系。也正因如此，在每次优化迭代中通常需要额外求解一次伴随方程以获得精确梯度信息。

接下来计算目标泛函对密度场的导数。根据残差算子的定义：
\[
\mathcal{R}(\rho,\boldsymbol{u})(\boldsymbol{\lambda}) = a_\rho(\boldsymbol{u},\boldsymbol{\lambda}) + s(\boldsymbol{u},\boldsymbol{\lambda}) - \ell_{\mathrm{in}}(\boldsymbol{\lambda}),
\]
可得其关于设计变量 $\rho$ 的偏 Gâteaux 导数为
\[
\frac{\partial\mathcal{R}(\rho,\boldsymbol{u}_\rho)(\boldsymbol{\lambda})}{\partial\rho}[\delta\rho] = \frac{\partial a_\rho(\boldsymbol{u}_\rho,\boldsymbol{\lambda})}{\partial\rho}[\delta\rho] + \frac{\partial s(\boldsymbol{u}_\rho,\boldsymbol{\lambda})}{\partial\rho}[\delta\rho] - \frac{\partial\ell_{\mathrm{in}}(\boldsymbol{\lambda})}{\partial\rho}[\delta\rho].
\]
鉴于输入/输出弹簧刚度 $k_{\mathrm{in}}, k_{\mathrm{out}}$、特征方向 $\boldsymbol{d}_{\mathrm{in}}, \boldsymbol{d}_{\mathrm{out}}$ 以及局部边界 $\Gamma_{\mathrm{in}}, \Gamma_{\mathrm{out}}$ 均为纯粹的物理设定，双线性型 $s(\cdot, \cdot)$ 与线性泛函 $\ell_{\mathrm{in}}(\cdot)$ 对设计变量 $\rho$ 均无显式依赖，即后两项的偏导数为零。因此上式简化为
\[
\frac{\partial\mathcal{R}(\rho,\boldsymbol{u}_\rho)(\boldsymbol{\lambda})}{\partial\rho}[\delta\rho] = \frac{\partial a_\rho(\boldsymbol{u}_\rho,\boldsymbol{\lambda})}{\partial\rho}[\delta\rho],
\]
另一方面，由双线性型 $a_\rho(\cdot,\cdot)$ 的定义：
\[
a_\rho(\boldsymbol{u},\boldsymbol{v}) = \int_{\Omega}\big(\mathbb{C}(\rho):\boldsymbol{\varepsilon}(\boldsymbol{u})\big) :\boldsymbol{\varepsilon}(\boldsymbol{v})\,\mathrm{d}\boldsymbol{x},
\]
并结合~\eqref{eq:ch2_local_linear_relationship} 所给出的局部线性关系，可得
\[
\begin{aligned} 
	\frac{\partial a_\rho(\boldsymbol{u}_\rho,\boldsymbol{\lambda})}{\partial\rho}[\delta\rho] &= \int_{\Omega} \big(\mathbb{C}'(\rho)[\delta\rho]:\boldsymbol{\varepsilon}(\boldsymbol{u}_\rho)\big) :\boldsymbol{\varepsilon}(\boldsymbol{\lambda})\,\mathrm{d}\boldsymbol{x}\\ &= \int_{\Omega} \big(\mathbb{C}'(\rho(\boldsymbol{x})):\boldsymbol{\varepsilon}(\boldsymbol{u}_\rho(\boldsymbol{x}))\big) : \boldsymbol{\varepsilon}(\boldsymbol{\lambda}(\boldsymbol{x}))\, \delta\rho(\boldsymbol{x})\,\mathrm{d}\boldsymbol{x}, 
\end{aligned}
\]
当伴随变量 $\boldsymbol{\lambda}$ 严格满足伴随方程时，由于未知状态扰动已被完全消除，因此输出位移约化目标泛函 $u_{\mathrm{out}}(\rho)$ 沿方向 $\delta\rho$ 的一阶总变分等于拉格朗日泛函对密度的偏导数：
\[
\delta u_{\mathrm{out}}(\rho) = \delta\mathcal{L}(\rho,\boldsymbol{u}_\rho,\boldsymbol{\lambda}) = \frac{\partial\mathcal{L}(\rho,\boldsymbol{u}_\rho,\boldsymbol{\lambda})}{\partial\rho}[\delta\rho] = \frac{\partial\mathcal{R}(\rho,\boldsymbol{u}_\rho)(\boldsymbol{\lambda})}{\partial\rho}[\delta\rho].
\]
从而得到变分的显式积分表达
\[
\delta u_{\mathrm{out}}(\rho) = \int_{\Omega} \big(\mathbb{C}'(\rho):\boldsymbol{\varepsilon}(\boldsymbol{u}_\rho)\big) :\boldsymbol{\varepsilon}(\boldsymbol{\lambda})\, \delta\rho\,\mathrm{d}\boldsymbol{x}.
\]
将其与变分的标准形式
\[
\delta u_{\mathrm{out}}(\rho) = \int_{\Omega}\frac{\delta u_{\mathrm{out}}}{\delta\rho}(\boldsymbol{x})\, \delta\rho(\boldsymbol{x})\,\mathrm{d}\boldsymbol{x}
\]
进行对比，即可直接提取出柔顺机构输出位移关于密度场的灵敏度密度函数：
\begin{equation} \label{eq:ch2_mechanism_sensitivity_general}
	\frac{\delta u_{\mathrm{out}}}{\delta\rho}(\boldsymbol{x}) = \big(\mathbb{C}'(\rho(\boldsymbol{x})) : \boldsymbol{\varepsilon}(\boldsymbol{u}_\rho(\boldsymbol{x}))\big) : \boldsymbol{\varepsilon}(\boldsymbol{\lambda}(\boldsymbol{x})).
\end{equation}
若采用经典的 SIMP 材料插值模型，该灵敏度密度可具体写为
\[
\frac{\delta u_{\mathrm{out}}}{\delta\rho}(\boldsymbol{x}) = E'(\rho(\boldsymbol{x}))\, \big(\mathbb{C}_0:\boldsymbol{\varepsilon}(\boldsymbol{u}_\rho(\boldsymbol{x}))\big) : \boldsymbol{\varepsilon}(\boldsymbol{\lambda}(\boldsymbol{x})).
\]


\subsubsection*{体积分数约束函数的灵敏度}

由第~\ref{subsec:min_compliance} 节可知，设设计域的总体积为 $V_0 = \int_{\Omega} 1\,\mathrm{d}\boldsymbol{x}$，材料的实际占用体积为 $V(\rho) = \int_{\Omega}\rho(\boldsymbol{x})\,\mathrm{d}\boldsymbol{x}$，则体积分数约束泛函定义为	
\[
g_V(\rho) := \frac{V(\rho)}{V_0} - V_f \le 0.
\]
由于 $g_V(\rho)$ 仅通过线性积分泛函显式依赖于设计变量 $\rho$，且与系统的状态位移场 $\boldsymbol{u}_\rho$ 无关，因此其灵敏度可以直接由 Gâteaux 导数给出，而无需构造拉格朗日泛函或引入伴随场。

对于任意容许的扰动方向 $\delta\rho$，体积分数约束的一阶变分可直接写为
\[
\delta g_V(\rho) = \frac{1}{V_0}\int_{\Omega}\delta\rho(\boldsymbol{x})\,\mathrm{d}\boldsymbol{x},
\]
将其与变分的标准积分形式
\[
\delta g_V(\rho) = \int_{\Omega}\frac{\delta g_V}{\delta\rho}(\boldsymbol{x})\,\delta\rho(\boldsymbol{x})\,\mathrm{d}\boldsymbol{x},
\]
进行对比，即可直接得到体积分数约束关于设计变量的灵敏度密度函数。该函数在整个设计域内为一常数：
\begin{equation} \label{eq:ch2_volume_sensitivity}
	\frac{\delta g_V}{\delta\rho}(\boldsymbol{x}) = \frac{1}{V_0}, \qquad \forall\,\boldsymbol{x}\in\Omega.
\end{equation}


\subsubsection*{局部应力约束函数的灵敏度}

由第~\ref{subsec:stress_constrained} 节可知，局部应力约束采用 von Mises 等效应力进行评估：
\[
\sigma^{\mathrm{vM}}(\rho, \boldsymbol{u}_\rho)(\boldsymbol{x}) \le \bar{\sigma}, \qquad  \forall \boldsymbol{x}\in \Omega\setminus\Omega_{\mathrm{void}}.
\]
为便于灵敏度分析，首先将位移场 $\boldsymbol{u}$ 与设计变量 $\rho$ 视为独立的自变量，定义未约化的点态约束泛函（归一化形式）：
\begin{equation}
	g_\sigma(\rho, \boldsymbol{u}; \boldsymbol{x}) := \frac{\sigma^{\mathrm{vM}}(\rho, \boldsymbol{u})(\boldsymbol{x})}{\bar{\sigma}} - 1 \le 0, \qquad \boldsymbol{x}\in \Omega\setminus\Omega_{\mathrm{void}}.
	\label{eq:unreduced_constraint}
\end{equation}
将满足平衡状态的位移解 $\boldsymbol{u} = \boldsymbol{u}_\rho$ 代入 \eqref{eq:unreduced_constraint} 式后，即可得到仅依赖于设计变量的约化约束泛函
\[
\tilde g_\sigma(\rho;\boldsymbol{x}) := g_\sigma(\rho,\boldsymbol{u}_\rho;\boldsymbol{x}),
\]
与前述的体积分数约束截然不同，约化应力约束泛函 $\tilde{g}_\sigma(\rho; \boldsymbol{x})$ 既通过材料插值 $\mathbb{C}(\rho)$ 对应力评估产生显式的局部依赖，又通过状态方程的耦合对位移场 $\boldsymbol{u}_\rho$ 产生高度非线性的隐式依赖。需要指出的是，上述点态应力约束及其灵敏度推导主要用于刻画局部应力约束的连续形式伴随结构。从严格的函数分析角度看，点值应力一般并不构成 $H^1(\Omega;\mathbb{R}^d)$ 空间上的连续泛函。因此，下述推导应理解为形式推导。在实际数值实现中，局部应力通常需在离散评估点上定义，并结合应力松弛或约束聚合等策略加以处理。

由链式法则，约化约束泛函 $\tilde{g}_\sigma(\rho;\boldsymbol{x})$ 沿方向 $\delta\rho$ 的一阶变分可展开为
\[
\delta \tilde g_\sigma(\rho;\boldsymbol{x}) = \frac{\partial g_\sigma}{\partial \rho}(\rho,\boldsymbol{u}_\rho;\boldsymbol{x})[\delta\rho] + \frac{\partial g_\sigma}{\partial \boldsymbol{u}}(\rho,\boldsymbol{u}_\rho;\boldsymbol{x})[\delta\boldsymbol{u}],
\]
其中，未知状态扰动 $\delta\boldsymbol{u}$ 由状态方程的线性化确定。由于状态解 $\boldsymbol{u}_\rho$ 满足
\[
a_\rho(\boldsymbol{u}_\rho,\boldsymbol{v})=\ell(\boldsymbol{v}),\qquad \forall \boldsymbol{v}\in \mathcal{V}_0,
\]
对其两端关于设计变量 $\rho$ 求变分可得
\[
a_\rho(\delta\boldsymbol{u},\boldsymbol{v}) = -\frac{\partial a_\rho(\boldsymbol{u}_\rho,\boldsymbol{v})}{\partial \rho}[\delta\rho], \qquad \forall \boldsymbol{v}\in \mathcal{V}_0.
\]
为避免显式求解全域的大规模扰动场 $\delta\boldsymbol{u}$，针对给定的应力评估点 $\boldsymbol{x}$，引入伴随变量 $\boldsymbol{\lambda}_\sigma^{(\boldsymbol{x})}\in \mathcal{V}_0$，并令其满足如下伴随方程：
\begin{equation} 
	a_\rho(\boldsymbol{v},\boldsymbol{\lambda}_\sigma^{(\boldsymbol{x})}) = -\frac{\partial g_\sigma}{\partial \boldsymbol{u}}(\rho,\boldsymbol{u}_\rho;\boldsymbol{x})[\boldsymbol{v}], \qquad \forall \boldsymbol{v}\in \mathcal{V}_0.
	\label{eq:ch2_stress_adjoint_equation}
\end{equation}
将 $\delta\boldsymbol{u}$ 作为测试函数代入 \eqref{eq:ch2_stress_adjoint_equation} 式，并结合状态方程的变分形式，即可消去 $\delta\boldsymbol{u}$ 项，得到点态应力约束的标准伴随灵敏度表达式：
\begin{equation} \label{eq:ch2_stress_sensitivity_standard}
	\delta \tilde{g}_\sigma(\rho;\boldsymbol{x}) = \frac{\partial g_\sigma}{\partial \rho}(\rho,\boldsymbol{u}_\rho;\boldsymbol{x})[\delta\rho] - \frac{\partial a_\rho(\boldsymbol{u}_\rho,\boldsymbol{\lambda}_\sigma^{(\boldsymbol{x})})}{\partial \rho}[\delta\rho].
\end{equation}
引入空间积分哑变量 $\boldsymbol{y} \in \Omega$，上式中的第二项（即双线性型对设计变量的偏导数）可进一步展开为
\[
\frac{\partial a_\rho(\boldsymbol{u}_\rho,\boldsymbol{\lambda}_\sigma^{(\boldsymbol{x})})}{\partial \rho}[\delta\rho] = \int_{\Omega} \big(\mathbb{C}'(\rho(\boldsymbol{y}))\,\delta\rho(\boldsymbol{y}) : \boldsymbol{\varepsilon}(\boldsymbol{u}_\rho(\boldsymbol{y}))\big) : \boldsymbol{\varepsilon}(\boldsymbol{\lambda}_\sigma^{(\boldsymbol{x})}(\boldsymbol{y}))\; \mathrm{d}\boldsymbol{y}.
\]

下面利用 Voigt 记号，给出伴随方程右端载荷项 $\partial g_\sigma / \partial \boldsymbol{u}$ 与显式偏导数项 $\partial g_\sigma / \partial \rho$ 的具体计算表达。根据 von Mises 等效应力的二次型定义：
\[
\sigma^{\mathrm{vM}}(\rho,\boldsymbol{u})(\boldsymbol{x}) = \sqrt{ \hat{\boldsymbol{\sigma}}(\rho,\boldsymbol{u})(\boldsymbol{x})^{\top}\boldsymbol{M}\, \hat{\boldsymbol{\sigma}}(\rho,\boldsymbol{u})(\boldsymbol{x}) },
\]
利用权重矩阵 $\boldsymbol{M}$ 的对称性，对任意应力场扰动 $\delta\hat{\boldsymbol{\sigma}}$，其一阶变分为
\[
\delta\sigma^{\mathrm{vM}}(\boldsymbol{x}) = \left(\frac{\boldsymbol{M}\hat{\boldsymbol{\sigma}}(\boldsymbol{x})}{\sigma^{\mathrm{vM}}(\boldsymbol{x})}\right)^{\!\top}\delta\hat{\boldsymbol{\sigma}}(\boldsymbol{x}).
\]

在代入状态解 $\boldsymbol{u} = \boldsymbol{u}_\rho$ 的前提下，当密度场 $\rho$ 固定时，应力向量关于位移场沿测试函数 $\boldsymbol{v}$ 方向的偏变分为 $\delta\hat{\boldsymbol{\sigma}}(\boldsymbol{x}) = \hat{\mathbb{C}}(\rho(\boldsymbol{x}))\,\hat{\boldsymbol{\varepsilon}}(\boldsymbol{v})(\boldsymbol{x})$。将其代入链式法则，可得到伴随方程中对应的虚拟载荷泛函项为：
\begin{equation} \label{eq:ch2_stress_virtual_load}
	\frac{\partial g_\sigma}{\partial \boldsymbol{u}}(\rho,\boldsymbol{u}_\rho;\boldsymbol{x})[\boldsymbol{v}] = \frac{1}{\bar{\sigma}} \left(\frac{\boldsymbol{M}\hat{\boldsymbol{\sigma}}(\boldsymbol{u}_\rho)(\boldsymbol{x})}{\sigma^{\mathrm{vM}}(\boldsymbol{u}_\rho)(\boldsymbol{x})}\right)^{\!\top} \Big(\hat{\mathbb{C}}(\rho(\boldsymbol{x}))\,\hat{\boldsymbol{\varepsilon}}(\boldsymbol{v})(\boldsymbol{x})\Big).
\end{equation}
同理，当位移场 $\boldsymbol{u}_\rho$ 固定时，应力向量关于密度场沿方向 $\delta\rho$ 的显式偏变分为 $\delta\hat{\boldsymbol{\sigma}}(\boldsymbol{x}) = \big(\hat{\mathbb{C}}'(\rho(\boldsymbol{x}))\,\delta\rho(\boldsymbol{x})\big)\, \hat{\boldsymbol{\varepsilon}}(\boldsymbol{u}_\rho)(\boldsymbol{x})$。由此可得应力约束对密度的显式偏导数项为：
\begin{equation} \label{eq:ch2_stress_explicit_derivative}
	\frac{\partial g_\sigma}{\partial \rho}(\rho,\boldsymbol{u}_\rho;\boldsymbol{x})[\delta\rho] = \frac{1}{\bar{\sigma}} \left(\frac{\boldsymbol{M}\hat{\boldsymbol{\sigma}}(\boldsymbol{u}_\rho)(\boldsymbol{x})}{\sigma^{\mathrm{vM}}(\boldsymbol{u}_\rho)(\boldsymbol{x})}\right)^{\!\top} \Big(\big(\hat{\mathbb{C}}'(\rho(\boldsymbol{x}))\,\delta\rho(\boldsymbol{x})\big)\, \hat{\boldsymbol{\varepsilon}}(\boldsymbol{u}_\rho)(\boldsymbol{x})\Big).
\end{equation}

综合式~\eqref{eq:ch2_stress_virtual_load}、\eqref{eq:ch2_stress_explicit_derivative} 与式 \eqref{eq:ch2_stress_sensitivity_standard}，可得点态应力约束 $\tilde g_\sigma(\rho;\boldsymbol{x})$ 沿任意设计扰动方向 $\delta \rho$ 的一阶变分为：
\begin{equation} \label{eq:ch2_stress_total_variation}
	\begin{aligned} 
		\delta \tilde g_\sigma(\rho;\boldsymbol{x}) ={}& \frac{1}{\bar{\sigma}} \left( \frac{ \boldsymbol{M}\hat{\boldsymbol{\sigma}}(\boldsymbol{u}_\rho)(\boldsymbol{x}) }{ \sigma^{\mathrm{vM}}(\boldsymbol{u}_\rho)(\boldsymbol{x}) } \right)^\top \Big( \big( \hat{\mathbb{C}}'(\rho(\boldsymbol{x}))\,\delta \rho(\boldsymbol{x}) \big) \hat{\boldsymbol{\varepsilon}}(\boldsymbol{u}_\rho)(\boldsymbol{x}) \Big) \\ &\quad - \int_\Omega \big( \mathbb{C}'(\rho(\boldsymbol{y}))\,\delta \rho(\boldsymbol{y}) : \boldsymbol{\varepsilon}(\boldsymbol{u}_\rho(\boldsymbol{y})) \big) : \boldsymbol{\varepsilon}\big(\boldsymbol{\lambda}_\sigma^{(\boldsymbol{x})}(\boldsymbol{y})\big) \,\mathrm{d}\boldsymbol{y}. 
	\end{aligned}
\end{equation}
该式表明，点态应力约束的灵敏度由显式项和伴随项两部分组成：前者反映材料插值对应力评估的直接影响，后者反映状态方程耦合所引起的隐式位移效应。


\section{拓扑优化中的数值优化算法}
\label{sec:num_algorithms}

拓扑优化问题的求解算法主要有优化准则（Optimality Criteria, OC）法和数学规划（Mathematical Programming, MP）法~\cite{GongChengJieGouYouHuaSheJiJiChuChengGengDongBianZhu2012}。依据优化迭代过程中对灵敏度信息的依赖程度，数学规划法通常被划分为无梯度算法与基于梯度的算法两大类。无梯度算法（如单纯形法、Powell 法及各类直接搜索法）虽然规避了复杂的灵敏度求解过程，但在面对拓扑优化通常涉及的数以万计甚至百万级设计变量时，其搜索效率显著下降，难以满足大规模问题的计算需求。

相比之下，基于梯度的优化算法通过利用目标函数与约束函数的导数信息来构造搜索方向或近似子问题，具有更高的求解效率和更好的可扩展性，因此成为连续体结构拓扑优化领域的主流选择。该类方法的典型代表包括序列线性规划（Sequential Linear Programming, SLP）、序列二次规划（Sequential Quadratic Programming, SQP）、序列凸规划（Sequential Convex Programming, SCP）以及专门针对结构优化特点发展的移动渐近线法（Method of Moving Asymptotes，MMA）等~\cite{qianApproachStructuralOptimization1984, svanbergMethodMovingAsymptotes1987, boggsSequentialQuadraticProgramming1995}。

本节主要介绍拓扑优化中常用的两类数值优化算法，OC 法与 MMA，并分别说明其基本思想、迭代格式及其在拓扑优化中的应用。


\subsection{优化准则法}
\label{subsec:alg_oc}

OC 法是一类基于极值条件显式更新设计变量的迭代算法~\cite{saveStructuralOptimizationVolume1986}。其基本思想是：针对给定的目标函数与约束条件，依据最优解所必须满足的 Karush–Kuhn–Tucker（KKT）条件推导出设计变量的最优性准则，并据此构造具有启发式特征的迭代更新格式。对于约束数目较少、结构较规整的拓扑优化问题（如单一体积分数约束下的柔顺度最小化问题），OC 法因其物理意义直观、实现简单且迭代代价较低，得到了广泛应用。

以单一体积分数约束问题为例，记离散后的设计变量向量为 $\boldsymbol{\rho} = (\rho_1,\dots,\rho_{N_d})^\top$，且满足基本的箱型约束 
\[
\rho_{\min} \le \rho_i \le 1,\quad{i}=1,\cdots,N_d.
\]
引入非负拉格朗日乘子 $\lambda$，针对目标函数 $c(\boldsymbol{\rho})$ 与体积约束 $g_V(\boldsymbol{\rho}) \le 0$ 构造简化的拉格朗日函数
\[
\mathcal{L}(\boldsymbol{\rho},\lambda) = c(\boldsymbol{\rho}) + \lambda{g}_V(\boldsymbol{\rho}).
\]
当体积约束处于活跃状态，且第 $i$ 个设计变量未触及上下界，即满足 $\rho_{\min} < \rho_i < 1$，则对应的 KKT 驻点条件要求 $\partial\mathcal{L}/\partial\rho_i = 0$。据此可定义第 $i$ 个变量的局部更新系数：
\[
B_i = -\frac{\partial{c}(\boldsymbol{\rho})}{\partial\rho_i}\left(\lambda\frac{\partial{g}_V(\boldsymbol{\rho})}{\partial\rho_i}\right)^{-1},
\]
在最优解处，内部变量原则上应满足 $B_i = 1$。这一关系反映了目标函数局部下降趋势与约束代价之间的平衡，是 OC 法构造显式更新格式的基本依据。

在实际迭代中，为同时满足箱型约束并控制单步更新幅度，通常在上述最优性关系的基础上进一步引入移动限制与阻尼机制。Bendsøe~\cite{bendsoeOptimizationStructuralTopology1995a} 提出了经典的 OC 更新格式：
\begin{equation} 
	\rho_i^{(k+1)} = 
	\begin{cases} 
		\max(\rho_{\min},\rho_i^{(k)}-m),\quad&\text{如果}\quad \rho_i^{(k)}(B_i^{(k)})^\eta \leq \max(\rho_{\min},\rho_i^{(k)}-m),\\ 
		\min(1,\rho_i^{(k)}+m),\quad&\text{如果}\quad \rho_i^{(k)}(B_i^{(k)})^\eta \geq \min(1,\rho_i^{(k)}+m),\\ 
		\rho_i^{(k)}(B_i^{(k)})^\eta,\quad&\text{其它情况}, 
	\end{cases}
	\label{eq:ch2_oc_update}
\end{equation}
其中，$m$ 为移动限制，用于约束单步迭代中设计变量的最大更新幅度，以防非线性迭代发散；$\eta$ 为阻尼系数，用于平滑收敛轨迹并抑制数值振荡。更新格式 ~\eqref{eq:ch2_oc_update} 在严格保证 $\rho_i \in [\rho_{\min}, 1]$ 的同时，通过系数 $(B_i^{(k)})^\eta$ 的自适应调节，驱动材料向高效承载区域集中，并从低效区域剔除。这里显式引入的极小正数 $\rho_{\min} > 0$ 作为密度下界，可有效避免有限元刚度矩阵因局部退化为孔洞而陷入奇异状态，从而保证了物理场求解的数值适定性。

基于上述密度更新公式，算法~\ref{alg:ch2_oc} 给出了优化准则法求解单一体积分数约束的拓扑优化问题的整体计算流程。该算法框架清晰地刻画了从物理场求解、灵敏度分析到密度更新与收敛判定的核心迭代闭环。
\begin{algorithm}[h]
	\caption{OC 法求解体积分数约束的拓扑优化问题}
	\label{alg:ch2_oc}
	\begin{spacing}{1.3}
		\begin{algorithmic}[1]
			\Require 初始设计变量 $\boldsymbol{\rho}^{(0)}$（满足
			$\rho_{\min} \leq \rho_i^{(0)} \leq 1$），移动限制 $m$，
			阻尼系数 $\eta$，收敛容差 $\epsilon$，最大迭代次数 $N_{\max}$
			\Ensure 满足最优性条件的设计变量 $\boldsymbol{\rho}^*$
			\State 令迭代计数器 $k \gets 0$，初始误差 $e \gets \infty$
			\While{$e \geq \epsilon$ 且 $k < N_{\max}$}
			\State 由 $\boldsymbol{\rho}^{(k)}$ 组装刚度矩阵，求解平衡方程，
			得位移场 $\boldsymbol{U}^{(k)}$
			\State 计算目标函数 $c(\boldsymbol{\rho}^{(k)})$
			及约束函数 $g_V(\boldsymbol{\rho}^{(k)})$
			\State 计算灵敏度 $\nabla c(\boldsymbol{\rho}^{(k)})$
			与 $\nabla g_V(\boldsymbol{\rho}^{(k)})$
			\State \textit{（可选）} 对灵敏度施加灵敏度过滤
			\State 二分法求拉格朗日乘子 $\lambda^* \geq 0$，使经 OC 格式
			更新的 $\boldsymbol{\rho}^{\mathrm{new}}$ 满足
			$g_V(\boldsymbol{\rho}^{\mathrm{new}}) \leq 0$
			\State \textit{（可选）} 对 $\boldsymbol{\rho}^{\mathrm{new}}$
			施加密度过滤或投影
			\State 计算当前更新误差
			$e \gets \|\boldsymbol{\rho}^{\mathrm{new}}
			- \boldsymbol{\rho}^{(k)}\|_{\infty,\Omega}$
			\State 令
			$\boldsymbol{\rho}^{(k+1)} \gets \boldsymbol{\rho}^{\mathrm{new}}$，并更新迭代计数器
			$k \gets k + 1$
			\EndWhile
			\State 返回 $\boldsymbol{\rho}^* \gets \boldsymbol{\rho}^{(k)}$
		\end{algorithmic}
	\end{spacing}
\end{algorithm}

尽管 OC 法在求解单一体积约束等常规问题时极为高效，但其主要局限性在于难以处理多约束与高度非线性的复杂拓扑优化模型。当面临多个不等式约束时，求解耦合的拉格朗日乘子会使其完全丧失显式更新的计算优势；而在处理柔顺机构设计或局部应力约束等非自伴随与非线性问题时，其启发式更新格式极易引发数值震荡。因此，针对多约束或复杂物理响应的拓扑优化问题，需要引入具有严格凸近似理论支撑的 MMA 以保障算法的全局鲁棒性。


\subsection{移动渐近线方法}
\label{subsec:alg_mma}

MMA 由 Svanberg~\cite{svanbergMethodMovingAsymptotes1987} 提出，是 SCP 方法的经典代表。在面临多约束和高度非线性的复杂拓扑优化问题时，MMA 展现出了极佳的适用性与鲁棒性。其核心数学思想是：在每次优化迭代中，通过引入一组动态更新的移动渐近线参数，在当前设计点附近构造目标与约束函数的可分严格凸近似。由此，原始的高维非线性优化问题被转化为一系列显式的严格凸子问题。得益于子问题具有的严格凸性与变量可分性，其既可以在原变量空间中直接求解，也可以借助拉格朗日对偶理论在低维对偶空间中高效求解，从而获取新的设计变量。随着渐近线参数的自适应调整，凸子问题序列将逐步逼近原问题，当满足收敛判据时，迭代解即可收敛至原非线性规划问题的 KKT 点。

一般形式的优化问题的数学模型可表示为：
\[
\begin{aligned}
	\min_{\boldsymbol{\rho}}\quad&f_0(\boldsymbol{\rho})\\
	\text{subject~to}\quad&f_i(\boldsymbol{\rho}) \leq 0,\quad&{i}=1,\cdots,N_g,\\
	\quad&\rho_j^{\min}\leq{\rho}_j\leq{\rho}_j^{\max},\quad&{j}=1,\cdots,N_d,
\end{aligned}
\]
其中，$\boldsymbol{\rho}$ 是设计变量向量，$f_0(\boldsymbol{\rho})$ 表示目标函数，$f_i(\boldsymbol{\rho}) \leq 0$ 为约束函数，$\rho_j^{\min}$ 和 $\rho_j^{\max}$ 分别为第 $j$ 个设计变量的下限和上限，$N_g$ 是约束函数的总个数。

在第 $k$ 次迭代中，MMA 基于当前设计点 $\boldsymbol{\rho}^{(k)}$ 处的目标与各约束函数的函数值 $f_i(\boldsymbol{\rho}^{(k)})$ 及其一阶梯度 $\nabla f_i(\boldsymbol{\rho}^{(k)})$ （$i=0,1,\dots,N_g$），引入一组动态更新的渐近线参数 $L_j^{(k)}$ 和 $U_j^{(k)}$，构造如下严格凸的近似子问题：
\begin{equation} \label{eq:ch2_mma_subproblem}
	\begin{aligned}
		\min_{\boldsymbol{\rho},\boldsymbol{y}, z}\quad&\tilde{f}_0^{(k)}(\boldsymbol{\rho}) + a_0z + \sum_{i=1}^{N_g}\left(c_iy_i+\frac{1}{2}d_iy_i^2\right)\\
		\text{subject to}\quad&\tilde{f}_i^{(k)}(\boldsymbol{\rho}) - a_iz - y_i \leq 0, \quad i=1,\cdots,N_g,\\
		&\alpha_j^{(k)}\leq\rho_j\leq\beta_j^{(k)}, \quad j=1,\cdots,N_d,\\
		&y_i\geq 0, \, z\geq 0.
	\end{aligned}
\end{equation}
其中，$\boldsymbol{y}=(y_1,\dots,y_{N_g})^\top$ 和 $z$ 为非负的人工松弛变量，用以扩大子问题的可行域并增强数值求解的稳定性；$a_0,a_i,c_i,d_i$ 为给定的非负惩罚参数；$\alpha_j^{(k)},\beta_j^{(k)}$ 为第 $j$ 个设计变量在第 $k$ 次迭代中的局部移动界限，它们由原始变量界限与移动渐近线参数共同决定。

针对目标函数与各项约束函数（统一记为 $i = 0, 1, \dots, N_g$），其可分严格凸近似函数 $\tilde{f}_i^{(k)}(\boldsymbol{\rho})$ 的标准解析形式统一构造为：
\begin{equation} \label{eq:ch2_mma_approximation}
	\tilde{f}_i^{(k)}(\boldsymbol{\rho}) = \sum_{j=1}^{N_d}\left(\frac{p_{ij}^{(k)}}{U_{j}^{(k)}-\rho_j} + \frac{q_{ij}^{(k)}}{\rho_{j}-L_j^{(k)}}\right) + r_i^{(k)},
\end{equation}
其中，$L_j^{(k)}$ 和 $U_j^{(k)}$ 分别为第 $j$ 个设计变量在第 $k$ 次迭代中的移动渐近线下界和上界；系数 $p_{ij}^{(k)}, q_{ij}^{(k)}, r_i^{(k)}$ 由原函数 $f_i$ 在当前设计点 $\boldsymbol{\rho}^{(k)}$ 处的函数值和一阶梯度确定。其基本要求是保证近似函数在当前设计点与原函数保持一阶一致，即满足
\[
\tilde{f}_i^{(k)}(\boldsymbol{\rho}^{(k)}) = f_i(\boldsymbol{\rho}^{(k)}),\quad\nabla\tilde{f}_i^{(k)}(\boldsymbol{\rho}^{(k)}) = \nabla f_i(\boldsymbol{\rho}^{(k)}).
\]

近似函数的凸性与数值行为再很大程度上取决于移动渐近线 $L_j^{(k)}$ 和 $U_j^{(k)}$ 的位置。当 $k=0,1$ 时，由于缺乏足够的历史迭代信息，通常采用对称初始化策略：
\[
\begin{aligned} 
	L_j^{(k)} &= \rho_j^{(k)} - 0.5(\rho_j^{\max}-\rho_j^{\min}),\\ 
	U_j^{(k)} &= \rho_j^{(k)} + 0.5(\rho_j^{\max}-\rho_j^{\min}). 
\end{aligned}
\]
当 $k\ge 2$ 时，则可根据最近两步迭代中设计变量的演化趋势对渐近线位置进行自适应更新：
\[
\begin{aligned} 
	L_j^{(k)} &= \rho_j^{(k)} - \gamma_j^{(k)}(\rho_j^{(k-1)}-L_j^{(k-1)}),\\ 
	U_j^{(k)} &= \rho_j^{(k)} + \gamma_j^{(k)}(U_j^{(k-1)}-\rho_j^{(k-1)}).
\end{aligned}
\]
其中，缩放参数 $\gamma_j^{(k)}$ 的取值通常根据相邻两次迭代步中该变量更新量差值的符号进行自适应调整。在经典 MMA 实现中，其常用更新策略可写为如下分段形式：
\[
\gamma_j^{(k)} = 
\begin{cases} 
	0.7, & \text{若 } (\rho_j^{(k)}-\rho_j^{(k-1)})(\rho_j^{(k-1)}-\rho_j^{(k-2)}) < 0, \\ 
	1.2, & \text{若 } (\rho_j^{(k)}-\rho_j^{(k-1)})(\rho_j^{(k-1)}-\rho_j^{(k-2)}) > 0, \\ 
	1.0, & \text{若 } (\rho_j^{(k)}-\rho_j^{(k-1)})(\rho_j^{(k-1)}-\rho_j^{(k-2)}) = 0.
\end{cases}
\]
上述策略的核心在于通过追踪变量的演化轨迹，自适应地调整凸近似函数的局部曲率：当变量震荡时，取较小的 $\gamma_j^{(k)}$ 以收缩渐近线并增强稳定性；当变量单调变化时，取较大的 $\gamma_j^{(k)}$ 以适度扩张渐近线并提高收敛速度；当变量变化平缓时，则取中间值。这里的 0.7、1.0 和 1.2 属于经典 MMA 实现中的常用经验取值，并非理论上唯一固定的参数选择。此外，实际计算中通常还需对 $L_j^{(k)}$ 和 $U_j^{(k)}$ 作适当截断处理，使其与当前设计点保持合理距离，以避免数值奇异或近似失效。

为了避免子问题求解过程中设计变量过于接近渐近线而导致数值奇异（即分母趋近于零），同时限制单步更新幅度以保证近似的有效性，MMA 显式定义了第 $k$ 次迭代的局部移动界限 $\alpha_j^{(k)}$ 和 $\beta_j^{(k)}$。其选取规则如下：
\[
\begin{aligned} 
	\alpha_j^{(k)} &= \max\left\{ \rho_j^{\min},\; L_j^{(k)} + 0.1(\rho_j^{(k)} - L_j^{(k)}),\; \rho_j^{(k)} - 0.5(\rho_j^{\max} - \rho_j^{\min}) \right\},\\ 
	\beta_j^{(k)} &= \min\left\{ \rho_j^{\max},\; U_j^{(k)} - 0.1(U_j^{(k)} - \rho_j^{(k)}),\; \rho_j^{(k)} + 0.5(\rho_j^{\max} - \rho_j^{\min}) \right\}, 
\end{aligned}
\]
其中，0.1 和 0.5 也是经典 MMA 实现中的常用经验取值，其作用在于在原始变量界限、渐近线位置与单步更新幅度之间建立平衡。上述界限实际上隐含了三重约束：首先，设计变量必须位于原始物理界限 $[\rho_j^{\min}, \rho_j^{\max}]$ 内；其次，单步更新幅度不能超过总设计域的 50\%；最后，更新后的变量与渐近线之间必须保留至少 10\% 的相对距离，即满足
\[
-0.9(\rho_j^{(k)} - L_j^{(k)}) \leq \rho_j - \rho_j^{(k)} \leq 0.9(U_j^{(k)} - \rho_j^{(k)}),
\]
这一机制有效地保证了凸近似子问题在数值求解时的非奇异性和稳定性。

为了满足前述的一阶一致性条件，MMA 算法对目标与约束函数在当前迭代点 $\boldsymbol{\rho}^{(k)}$ 处的局部偏导数进行了正负部分离。记第 $i$ 个函数关于第 $j$ 个设计变量偏导数的正部与负部分别为
\[
\left(\frac{\partial f_i}{\partial \rho_j}\right)^{+} = \max\left(\frac{\partial f_i}{\partial \rho_j}(\boldsymbol{\rho}^{(k)}), 0\right), \quad  \left(\frac{\partial f_i}{\partial \rho_j}\right)^{-} = \max\left(-\frac{\partial f_i}{\partial \rho_j}(\boldsymbol{\rho}^{(k)}), 0\right),
\]
基于此，可进一步构造近似函数所需的非负系数：
\[
\begin{aligned} 
	p_{ij}^{(k)} &= (U_j^{(k)}-\rho_j^{(k)})^2 \left[ 1.001\left(\frac{\partial f_i}{\partial \rho_j}\right)^{+} + 0.001\left(\frac{\partial f_i}{\partial \rho_j}\right)^{-} + \frac{10^{-5}}{\rho_j^{\max}-\rho_j^{\min}} \right],\\ 
	q_{ij}^{(k)} &= (\rho_j^{(k)}-L_j^{(k)})^2 \left[ 0.001\left(\frac{\partial f_i}{\partial \rho_j}\right)^{+} + 1.001\left(\frac{\partial f_i}{\partial \rho_j}\right)^{-} + \frac{10^{-5}}{\rho_j^{\max}-\rho_j^{\min}} \right], 
\end{aligned}
\]
其中，1.001、0.001 以及 $10^{-5}$ 均属于经典 MMA 实现中的常用经验参数。前两者通过对正、负梯度分量施加非对称权重，为凸近似引入适度的保守性，以增强非线性迭代的稳定性；后者则作为极小正数修正项，用于避免在局部梯度过小或退化情形下系数出现零值。上述构造有助于保持近似函数的严格凸性与数值稳定性，从而使子问题具有良好的可解性。在经典 MMA 参数设置下，子问题通常可获得唯一解。

综合上述关于移动渐近线更新、严格凸近似函数构造以及局部移动界限控制的理论推导，算法~\ref{alg:mma} 给出了基于 MMA 求解带约束拓扑优化问题的整体计算流程。该算法框架完整刻画了从物理场求解、灵敏度分析到近似子问题构造与解算的核心迭代闭环。
\begin{algorithm}[h]
	\caption{MMA 求解带约束的拓扑优化问题}
	\label{alg:mma}
	\begin{spacing}{1.3}
		\begin{algorithmic}[1]
			\Require 初始设计变量 $\boldsymbol{\rho}^{(0)}$（满足 $\rho_j^{\min} \leq \rho_j^{(0)} \leq \rho_j^{\max}$），MMA 惩罚参数 $a_0, a_i, c_i, d_i$，收敛容差 $\epsilon$，最大迭代次数 $N_{\max}$
			\Ensure 满足最优性条件的设计变量 $\boldsymbol{\rho}^*$
			\State 初始化：令迭代计数器 $k \gets 0$，初始误差 $e \gets \infty$，并给定初始渐近线 $L_j^{(0)}, U_j^{(0)}$
			\While{$e \geq \epsilon$ 且 $k < N_{\max}$}
			\State 由当前设计变量 $\boldsymbol{\rho}^{(k)}$ 求解状态方程，获取系统物理响应
			\State 计算目标函数 $f_0(\boldsymbol{\rho}^{(k)})$ 与各项约束函数 $f_i(\boldsymbol{\rho}^{(k)})$
			\State 计算目标函数的灵敏度 $\nabla f_0(\boldsymbol{\rho}^{(k)})$ 与约束函数的灵敏度 $\nabla f_i(\boldsymbol{\rho}^{(k)})$
			\State \textit{（可选）} 对灵敏度场施加灵敏度过滤
			\State 依据演化轨迹的历史信息，自适应更新第 $k$ 步的移动渐近线 $L_j^{(k)}$ 与 $U_j^{(k)}$
			\State 综合原始物理界限与当前渐近线位置，构造局部移动界限 $\alpha_j^{(k)}$ 与 $\beta_j^{(k)}$
			\State 计算非负系数 $p_{ij}^{(k)}, q_{ij}^{(k)}$，构造各目标与约束函数的可分严格凸近似 $\tilde{f}_i^{(k)}(\boldsymbol{\rho})$
			\State 在局部界限 $[\alpha_j^{(k)}, \beta_j^{(k)}]$ 内，采用原对偶内点算法求解 MMA 严格凸子问题，获取更新变量 $\boldsymbol{\rho}^{\mathrm{new}}$
			\State \textit{（可选）} 对 $\boldsymbol{\rho}^{\mathrm{new}}$ 施加密度过滤或投影
			\State 计算当前更新误差 $e \gets \|\boldsymbol{\rho}^{\mathrm{new}} - \boldsymbol{\rho}^{(k)}\|_{\infty,\Omega}$
			\State 令 $\boldsymbol{\rho}^{(k+1)} \gets \boldsymbol{\rho}^{\mathrm{new}}$，并更新迭代计数器 $k \gets k + 1$
			\EndWhile
			\State 返回 $\boldsymbol{\rho}^* \gets \boldsymbol{\rho}^{(k)}$
		\end{algorithmic}
	\end{spacing}
\end{algorithm}


\section{正则化与长度尺度控制：过滤与投影}
\label{sec:regularization}

连续体拓扑优化本质上是在无限维函数空间中寻求最优材料分布 $\rho(\boldsymbol{x})$。然而，由于缺乏必要的正则性约束，该问题在数学上通常是不适定的。在数值求解中，这直接表现为严重的网格依赖性与棋盘格现象——随着网格细化，结构会衍生出无限精细的微观分支，且无法收敛于确定的宏观构型。为了获得网格无关且具备可制造性的设计，必须引入正则化手段来控制结构中的最小长度尺度。目前，基于卷积积分的过滤技术及其衍生的非线性投影方法，构成了处理该问题最主流的数值计算框架之一~\cite{sigmundMorphologybasedBlackWhite2007b}。 


\subsection{灵敏度过滤方法}
\label{subsec:filter_sensitivity}

灵敏度过滤由 Sigmund~\cite{sigmundDesignCompliantMechanisms1997a} 提出。虽然它本质上是一种启发式方法，缺乏严格的数学变分基础，但其物理直观清晰，且能够有效消除棋盘格现象。在连续域 $\Omega$ 上，该方法通过对目标泛函的梯度场进行卷积平滑来修正搜索方向。

定义修正后的灵敏度场 $\widetilde{\frac{\delta\mathcal{J}}{\delta\rho}}(\boldsymbol{x})$ 为原始灵敏度场 $\frac{\delta\mathcal{J}}{\delta\rho}(\boldsymbol{y})$ 在局部邻域内的加权平均：
\begin{equation} \label{eq:ch2_sensitivity_filter}
	\widetilde{\frac{\delta\mathcal{J}}{\delta\rho}}(\boldsymbol{x}) = \frac{1}{\max\{\gamma, \rho(\boldsymbol{x})\}\psi(\boldsymbol{x})}\int_{\Omega}w(\|\boldsymbol{y} - \boldsymbol{x}\|) \rho(\boldsymbol{y}) \frac{\delta\mathcal{J}}{\delta\rho}(\boldsymbol{y}) \, \mathrm{d}\boldsymbol{y},
\end{equation}
式中，$\boldsymbol{x}, \boldsymbol{y} \in \Omega$ 为空间坐标向量；$\gamma$ 为防止奇异的极小正数；$\psi(\boldsymbol{x})$ 为归一化因子，定义为卷积核在设计域内的积分：
\[
\psi(\boldsymbol{x}) = \int_{\Omega}w(\Vert\boldsymbol{y}-\boldsymbol{x}\Vert)\,\mathrm{d}\boldsymbol{y},
\]
其中 $w(r)$ 为具有紧支集的卷积核函数，通常采用如下多项式衰减形式~\cite{brunsTopologyOptimizationNonlinear2001, bourdinFiltersTopologyOptimization2001}：
\[
w(r) = \max\{0, r_{\min} - r\}^q, \quad q \geq 1,
\]
其中 $r = \|\boldsymbol{y} - \boldsymbol{x}\|$ 表示两点间的欧几里得距离，$r_{\min}$ 为预设的过滤半径，$q$ 为控制衰减速率的指数。特别地，当 $q = 1$ 时，上式退化为常用的线性锥形核；当 $q > 1$ 时，权重随距离衰减更为迅速，邻域内远处点的影响被进一步削弱。上述积分形式表明，某空间点处的修正灵敏度不仅取决于该点本身，还受到其邻域内灵敏度场的加权影响，从而在数值计算层面平滑了灵敏度场的高频振荡。尽管灵敏度过滤在工程计算中应用广泛，但由于它并未直接改变设计变量场本身，因此其正则化作用主要体现在优化更新方向的修正上，而非对设计空间的严格约束。


\subsection{密度过滤方法}
\label{subsec:filter_density}

灵敏度过滤虽然在工程上有效，但其本质是对梯度场的后处理，并未改变设计变量本身的空间分布，因此在数学上缺乏对解空间的严格约束。相比之下，Bruns 和 Tortorelli~\cite{brunsTopologyOptimizationNonlinear2001} 提出的密度过滤直接作用于设计变量场，是一种更具理论基础的正则化手段。Bourdin~\cite{bourdinFiltersTopologyOptimization2001} 进一步证明了在适当条件下，该方法能够保证拓扑优化问题解的存在性。

密度过滤的核心思想在于解耦“数学设计空间”与“物理材料空间”。在该框架下，原始设计变量场 $\rho(\boldsymbol{x})$ 不再直接决定材料的物理属性，而是作为一组纯粹的数学优化参数。真正决定结构物理响应（如刚度、质量、热传导率等）的是经过平滑映射后的物理密度场 $\tilde{\rho}(\boldsymbol{x})$。因此，在仅采用密度过滤而尚未引入投影的情形下，状态方程求解及性能泛函计算均应基于 $\tilde{\rho}$ 进行，且最终的设计结果亦应以 $\tilde{\rho}$ 为准。

在连续域 $\Omega$ 上，密度过滤定义了一个从原始设计变量空间 $L^\infty(\Omega)$ 到物理密度空间的平滑映射算子。物理密度场 $\tilde{\rho}(\boldsymbol{x})$ 定义为原始设计场的卷积：
\begin{equation} \label{eq:ch2_density_filter}
	\tilde{\rho}(\boldsymbol{x}) = \frac{1}{\psi(\boldsymbol{x})}\int_{\Omega} w(\|\boldsymbol{y} - \boldsymbol{x}\|) \rho(\boldsymbol{y}) \, \mathrm{d}\boldsymbol{y},
\end{equation}
其中，$\psi(\boldsymbol{x})$ 与卷积核 $w(\cdot)$ 的定义与前述灵敏度过滤一致。由于卷积核满足非负性，且采用了归一化处理，若原始设计变量满足
\[
\rho(\boldsymbol{x})\in[\rho_{\min},1], \qquad \boldsymbol{x}\in\Omega,
\]
则过滤后的 $\tilde{\rho}(\boldsymbol{x})$ 仍保持在同一区间内。在此框架下，结构的目标泛函及约束泛函 $\mathcal{F}$ 均直接依赖于物理密度场 $\tilde{\rho}$，根据链式法则，其灵敏度可写为：
\begin{equation} \label{eq:ch2_density_filter_sensitivity}
	\frac{\delta\mathcal{F}}{\delta\rho(\boldsymbol{x})} = \int_{\Omega}\frac{\delta\mathcal{F}}{\delta\tilde{\rho}(\boldsymbol{y})}\frac{\delta\tilde{\rho}(\boldsymbol{y})}{\delta\rho(\boldsymbol{x})}\,\mathrm{d}\boldsymbol{y} = \int_{\Omega} w(\|\boldsymbol{y} - \boldsymbol{x}\|) \frac{1}{\psi(\boldsymbol{y})}\frac{\delta\mathcal{F}}{\delta\tilde{\rho}(\boldsymbol{y})}\, \mathrm{d}\boldsymbol{y}.
\end{equation}


\subsection{投影方法}
\label{subsec:projection_methods}

虽然密度过滤方法在数学上提供了良好的适定性证明，并有效控制了最小长度尺度，但其低通滤波特性不可避免地在结构边界处引入了较宽的中间密度过渡区域（即灰度带）。这导致最终优化结果中存在大量介于 $0$ 与 $1$ 之间的中间密度分布，不仅增加了制造难度，且在有限元分析中极易导致材料属性（如刚度）被低估或高估。为了获得边界清晰、黑白分明的拓扑结构，Guest 等~\cite{guestAchievingMinimumLength2004b} 及 Sigmund~\cite{sigmundMorphologybasedBlackWhite2007b} 提出在密度过滤之后引入基于 Heaviside 函数的非线性投影技术。

投影方法的核心思想是构建一个从平滑密度空间到物理密度空间的非线性映射。在该框架下，拓扑优化的变量体系被正式扩展为标准的 “三场表述”：
\[
\rho(\boldsymbol{x}) \xrightarrow{\text{过滤算子}} \tilde{\rho}(\boldsymbol{x}) \xrightarrow{\text{投影算子}} \bar{\rho}(\boldsymbol{x}),
\]
其中，$\rho(\boldsymbol{x})$ 为原始设计变量，$\tilde{\rho}(\boldsymbol{x})$ 为过滤后密度，$\bar{\rho}(\boldsymbol{x})$ 为最终用于状态方程求解和性能评估的物理密度。也就是说，在引入投影之后，$\tilde{\rho}$ 不再直接承担物理材料属性的角色，而是作为连接设计变量与最终物理密度之间的中间场。

理想的投影算子为 Heaviside 阶跃函数，可形式化表示为：
\[
\bar{\rho}(\boldsymbol{x}) = \mathcal{H}(\tilde{\rho}(\boldsymbol{x}) - \eta) = 
\begin{cases}
	1, & \text{若 } \tilde{\rho}(\boldsymbol{x}) > \eta, \\
	0, & \text{若 } \tilde{\rho}(\boldsymbol{x}) \leq \eta,
\end{cases}
\]
其中 $\eta \in (0,1)$ 为投影阈值。然而，标准的阶跃函数在阈值 $\eta$ 处不可微，因此无法直接应用于基于伴随方法和梯度的数值优化算法。为此，实际计算中通常采用光滑的连续函数来逼近理想的阶跃特征。目前常用的两类光滑投影函数分别为指数型投影与双曲正切型投影:
\begin{enumerate}
	\item \textbf{指数型投影（常被称为 Heaviside 投影）}：由 Guest 等~\cite{guestAchievingMinimumLength2004b} 提出，其形式为：
	\begin{equation} \label{eq:ch2_exponential_projection}
		\bar{\rho}(\boldsymbol{x}) = 1 - e^{-\beta \tilde{\rho}(\boldsymbol{x})} + \tilde{\rho}(\boldsymbol{x}) e^{-\beta},
	\end{equation}
	其中 $\beta > 0$ 为控制投影陡峭程度的正则化参数。当 $\beta \to 0$ 时，该映射退化为近似线性关系；当 $\beta$ 逐渐增大时，投影函数的非线性增强，输出结果趋于更加接近 0-1 分布。与显式引入阈值参数的投影形式相比，该方法不直接给出独立的阈值位置，更可视为一类单侧强化的光滑投影。在部分微结构设计与特征尺度控制问题中，该类投影也有应用。

	\item \textbf{双曲正切型投影（常被称为阈值投影）}：由 Wang 等~\cite{wangProjectionMethodsConvergence2011a} 推广，其表达式为：
	\begin{equation} \label{eq:ch2_tanh_projection}
		\bar{\rho}(\boldsymbol{x}) = \frac{\tanh(\beta \eta) + \tanh(\beta (\tilde{\rho}(\boldsymbol{x}) - \eta))}{\tanh(\beta \eta) + \tanh(\beta (1 - \eta))},
	\end{equation}
	式中，$\eta \in [0,1]$ 为投影阈值，用于控制灰度向黑白转换的分界位置，$\beta>0$ 仍控制投影的陡峭程度。该形式具有较好的对称性，且阈值参数可独立调节，因此在实际拓扑优化中更为常用。通常取 $\eta=0.5$，以在一定程度上保持过滤前后的体积分数近似守恒。
\end{enumerate}
投影参数 $\beta$ 的取值对优化收敛行为具有重要影响。若一开始即取较大的 $\beta$，则目标泛函往往表现出较强的非凸性，容易使优化过程过早陷入局部极小；若 $\beta$ 过小，则投影作用较弱，难以有效抑制灰度。为此，实际计算中常采用延拓策略：在优化初期取较小的 $\beta$ 值（如 $\beta=1$），此时优化问题近似凸性较好，利于寻找全局轮廓。随着迭代推进，再逐步增大 $\beta$，从而逐渐锐化结构边界并逼近清晰的黑白设计。

引入投影算子后，目标泛函 $\mathcal{F}$ 对原始设计变量 $\rho(\boldsymbol{x})$ 的灵敏度计算需根据链式法则进一步扩展。基于连续域的变分推导，其导数关系为
\[
\frac{\delta \mathcal{F}}{\delta \rho(\boldsymbol{x})} = \int_{\Omega} \frac{\delta \mathcal{F}}{\delta \bar{\rho}(\boldsymbol{y})} \frac{\mathrm{d} \bar{\rho}(\boldsymbol{y})}{\mathrm{d} \tilde{\rho}(\boldsymbol{y})} \frac{\delta \tilde{\rho}(\boldsymbol{y})}{\delta \rho(\boldsymbol{x})} \, \mathrm{d}\boldsymbol{y},
\]
考虑到投影是点对点的局部映射，即 $\bar{\rho}(\boldsymbol{y})$ 仅取决于该点的 $\tilde{\rho}(\boldsymbol{y})$，而密度过滤是全域积分映射，结合第 \ref{subsec:filter_density} 节的过滤伴随算子，可得最终灵敏度表达式：
\begin{equation} \label{eq:ch2_final_chain_rule}
	\frac{\delta \mathcal{F}}{\delta \rho(\boldsymbol{x})} = \int_{\Omega} w(\|\boldsymbol{y} - \boldsymbol{x}\|) \frac{1}{\psi(\boldsymbol{y})} \left( \frac{\delta \mathcal{F}}{\delta \bar{\rho}(\boldsymbol{y})} \cdot \frac{\mathrm{d} \bar{\rho}}{\mathrm{d} \tilde{\rho}} \right) \, \mathrm{d}\boldsymbol{y},
\end{equation}
其中 $\mathrm{d} \bar{\rho} / \mathrm{d} \tilde{\rho}$ 为局部投影函数的标量导数。对于指数型投影，其导数为
\[
\frac{\mathrm{d} \bar{\rho}}{\mathrm{d} \tilde{\rho}} = \beta e^{-\beta \tilde{\rho}} + e^{-\beta},
\]
而对于双曲正切型投影，其导数为
\[
\frac{\mathrm{d} \bar{\rho}}{\mathrm{d} \tilde{\rho}} = \beta \frac{1 - \tanh^2(\beta (\tilde{\rho} - \eta))}{\tanh(\beta \eta) + \tanh(\beta (1 - \eta))}.
\]
上述链式关系表明，物理场灵敏度在反向传播时，首先经由投影导数完成局部加权调制，随后再通过卷积核完成空间平滑传播。由于投影导数在灰度区域通常较大、而在接近纯黑白区域时相对较小，因此这种“过滤 + 投影”的组合机制能够将优化驱动力更集中地传递到结构边界附近，从而更有效地推动几何边界演化并提升最终拓扑构型的清晰度。


%\fi